\documentclass[11pt]{article}
\usepackage[utf8]{inputenc}
\usepackage{amsmath, amssymb, amsthm}
\usepackage{booktabs}
\usepackage{array}
\usepackage{graphicx}
\usepackage{hyperref}
\usepackage[numbers]{natbib}
\usepackage[margin=1in]{geometry}
\usepackage{tikz}

\usetikzlibrary{positioning,arrows.meta}
\newcommand{\keywords}[1]{\vspace{0.5em}\noindent\textbf{Keywords:} #1}

\title{Empirical-Frontier Regularization for LLM Synthetic Agents: A Proof-of-Concept Framework and Case Study in Vaccination Timing}
\author{Yichao Jin \\ University of Texas at Dallas \\ \texttt{Yichao.Jin@UTDallas.edu}}
\date{}




\begin{document}

\maketitle

\begin{abstract}

\textbf{Objective:}
Large language models (LLMs) are increasingly used as synthetic agents for healthcare decision support and policy simulation, yet existing approaches provide limited mechanisms for incorporating empirically estimated behavioral knowledge into inference. We propose \textbf{Empirical-Frontier Regularization (EFR)}, a transparent inference-time architecture that integrates behavioral knowledge derived from discrete choice experiments (DCEs) with LLM-generated decisions.

\textbf{Methods:}
EFR estimates an empirical choice frontier using a mixed-logit DCE model and combines this empirical behavioral reference with LLM-generated probabilities through a transparent convex anchoring rule. The framework operates entirely at inference time, requiring neither retraining nor fine-tuning while preserving explicit behavioral constraints on measured attributes. We further compare EFR with a uniform-shrinkage baseline to distinguish the effects of empirical regularization from retained LLM variation.

\textbf{Results:}
A proof-of-concept evaluation using a vaccination-timing DCE involving 1,027 adults demonstrated the feasibility of the proposed architecture. Across three open-source LLMs, EFR substantially improved behavioral consistency relative to unconstrained models, reducing waiting-time monotonicity violations from 62.5\% to 4.2\% and restoring rank-order agreement with the empirical behavioral reference ($\rho=-0.024$ to $0.952$). Comparison with the uniform-shrinkage baseline indicated that behavioral consistency was driven primarily by empirical regularization toward the DCE frontier, whereas retained LLM variation contributed limited additional signal within the original DCE design space.

\textbf{Conclusion:}
EFR demonstrates a transparent and auditable architecture for integrating empirically estimated behavioral knowledge with foundation models through inference-time regularization. Rather than replacing either DCE models or LLMs, EFR combines their complementary strengths: the empirical behavioral model provides the primary behavioral signal, while the LLM provides flexible representation for rich natural-language scenarios. More broadly, this work illustrates how empirically grounded statistical models can serve as transparent and auditable inference-time constraints for foundation models in biomedical informatics.

\noindent\textbf{Keywords:}
Large Language Models; Biomedical Informatics; Discrete Choice Experiments; Inference-Time Regularization; Behavioral Simulation; Clinical Decision Support.

\end{abstract}

\clearpage
\section{Introduction}
\label{sec:intro}

Large language models (LLMs) are increasingly used as synthetic agents for healthcare decision support, policy evaluation, and simulation~\citep{Park2023a, Mu2024, VonHoene2025}. Their principal advantage lies in their ability to interpret rich natural-language scenarios that extend beyond the structured attribute spaces supported by conventional statistical models. However, fluent language generation alone does not guarantee behaviorally consistent decision making. For example, in a vaccination policy simulation, an unconstrained LLM may assign a higher acceptance probability to a strategy with a longer waiting time than to an otherwise identical strategy with a shorter waiting time. Such reversals indicate that natural-language reasoning alone is insufficient for policy simulation, where generated decisions are expected to remain consistent with empirically observed behavioral relationships.

Recent studies have shown that LLMs can generate coherent synthetic respondents and simulate complex healthcare narratives~\citep{Aher2022TuringExperiments, Horton2023, santurkar2023whose}. However, existing approaches provide limited mechanisms for incorporating empirically estimated behavioral knowledge into the inference process. As a result, unconstrained LLMs may reverse attribute-specific preference ordering, making their simulated decisions difficult to interpret or trust in policy evaluation without empirical grounding. More fundamentally, existing approaches do not explicitly distinguish the behavioral information supplied by empirical data from the representational flexibility supplied by foundation models, making hybrid AI systems difficult to interpret and audit.

One source of empirically grounded behavioral knowledge is the discrete choice experiment (DCE), which estimates how observed human preferences vary with decision attributes~\citep{Chan2024}. In this study, DCEs are not treated as models to be replaced by LLMs, but as sources of empirically estimated behavioral structure. A DCE-estimated choice frontier specifies the sign and ordering of attribute-specific preferences supported by observed responses. The challenge is therefore not to replace structured choice models with foundation models, but to integrate empirical behavioral knowledge with the representational flexibility of LLMs through a transparent inference-time architecture.

To address this challenge, we propose \textbf{Empirical-Frontier Regularization (EFR)}, a transparent inference-time architecture for integrating empirically estimated behavioral knowledge into LLM-generated decisions. EFR estimates an empirical choice frontier from DCE responses and combines the LLM-generated probability with the empirical frontier through a transparent convex anchoring rule. The framework is intentionally designed to be transparent and auditable: the empirical behavioral reference, anchoring weight, raw LLM probability, and regularized output remain explicitly observable throughout the inference process. Rather than modifying the underlying language model, EFR performs probability-level regularization after generation, allowing the same foundation model to be grounded using different empirical datasets without retraining or fine-tuning.

We demonstrate the proposed architecture using a vaccination-timing DCE involving 1,027 adults in Wuhan, China. Rather than introducing a new predictive model, this study presents a proof-of-concept demonstration of transparent inference-time empirical regularization. To better understand the respective roles of empirical regularization and retained LLM variation, we additionally compare EFR with a uniform-shrinkage baseline that replaces the LLM probability with a global constant while preserving the same convex anchoring structure.

Specifically, this study makes three contributions:

\begin{enumerate}

\item \textbf{Methodological contribution.}
We introduce Empirical-Frontier Regularization (EFR), a lightweight and transparent inference-time architecture that injects empirically estimated behavioral knowledge into LLM-generated choice probabilities through a convex anchoring rule, requiring neither retraining nor fine-tuning.

\item \textbf{Empirical contribution.}
Using a real-world vaccination DCE, we evaluate the proposed architecture and compare it with a uniform-shrinkage baseline to disentangle the respective contributions of empirical regularization and retained LLM variation. This decomposition demonstrates that behavioral consistency within the original DCE design space is driven primarily by the empirical frontier rather than retained LLM variation.

\item \textbf{Conceptual contribution.}
We argue that empirical behavioral models and foundation models serve complementary roles within hybrid AI systems: empirical models provide transparent behavioral constraints, whereas LLMs provide flexible representation for rich natural-language policy scenarios. EFR operationalizes this role separation through a transparent and auditable inference-time architecture.

\end{enumerate}

More broadly, this study illustrates how empirically estimated statistical models can serve as transparent and auditable inference-time constraints for foundation models. Rather than positioning statistical models and generative AI as competing paradigms, EFR demonstrates a hybrid architecture in which empirical behavioral knowledge and language-model reasoning play complementary roles in trustworthy biomedical informatics systems.

The remainder of this study is organized as follows. Section~2 reviews related work. Section~3 presents the proposed framework. Section~4 reports the empirical evaluation. Section~5 discusses the methodological implications, limitations, and future directions of the proposed architecture.

\section{Related Work}

Our work connects two research areas: LLM-based synthetic agents for healthcare simulation and empirical behavioral models derived from discrete choice experiments (DCEs). While these two directions have developed largely independently, both seek to improve the realism and reliability of decision support. Our work lies at their intersection by incorporating empirically estimated behavioral knowledge into LLM-based synthetic agents through inference-time regularization.

\subsection{LLM-Based Synthetic Agents}

Large language models (LLMs) have enabled a new generation of synthetic agents capable of generating natural-language reasoning and simulating human decision making across diverse domains~\citep{Bommasani2021, Park2023a, Mu2024}. In healthcare, these agents have been explored for public health communication, policy evaluation, and behavioral simulation~\citep{Sehgal2025, Gullison2025}. Their principal advantage is representational flexibility: unlike conventional statistical models, LLMs can process rich textual descriptions and respond to scenarios that extend beyond predefined structured variables.

However, behavioral realism remains a major challenge. Although LLM-generated responses are often linguistically coherent, they may violate empirically established behavioral relationships or produce implausible preference orderings when evaluated against observed human decision patterns~\citep{Aher2022TuringExperiments, Horton2023, hintze2025nonlinear}. Existing approaches, including prompt engineering and reasoning strategies, can reduce some inconsistencies but provide no explicit mechanism for enforcing empirically grounded behavioral constraints throughout a policy space~\citep{wang2025simulate, wang2025heating}.

\subsection{Empirical Behavioral Models from Discrete Choice Experiments}

Discrete choice experiments (DCEs) provide a well-established framework for estimating how individuals trade off policy-relevant attributes such as waiting time, vaccine efficacy, side effects, and financial incentives~\citep{soekhai2019discrete, chan2024vaccine}. Under the random utility framework~\citep{train2009discrete, louviere2000stated}, DCE responses can be used to estimate an empirical choice frontier that preserves interpretable relationships between decision attributes and observed choices. Mixed-logit models further account for preference heterogeneity across respondents~\citep{hensher2003mixed, hensher2015applied}.

In this study, we treat the DCE frontier as an empirically estimated source of behavioral knowledge rather than as assumption-free ground truth. While empirical choice models provide transparent and behaviorally interpretable predictions within their original experimental design, they cannot directly operate on rich natural-language policy scenarios.

\subsection{Bridging Empirical Behavioral Models and LLM-Based Agents}

Despite these advances, an important methodological gap remains. Existing LLM-based synthetic agents provide representational flexibility but lack explicit mechanisms for enforcing empirically validated behavioral constraints\citep{li2026position, tornberg2025validation}. Conversely, empirical behavioral models derived from DCEs provide transparent and interpretable behavioral structure but remain confined to the structured experimental designs from which they are estimated. Consequently, existing approaches either prioritize flexible language understanding without behavioral guarantees or preserve behavioral validity at the cost of limited representational flexibility.

Recent studies have explored hybrid approaches that combine LLMs with structured behavioral or economic models to improve decision quality~\citep{sfeir2025assist, zhang2026agentic}. These methods demonstrate the value of incorporating external behavioral knowledge into LLM-based decision making, but they generally rely on prompt design, end-to-end learning, or implicit alignment mechanisms.

To our knowledge, no existing framework incorporates DCE-estimated empirical choice frontiers directly into LLM inference through a transparent inference-time regularization mechanism. Our work addresses this gap by treating empirically estimated behavioral knowledge as an explicit inference-time constraint rather than a target for parameter optimization or prompt engineering. In this way, our framework occupies a complementary position within the emerging landscape of knowledge-guided AI: it combines the representational flexibility of LLM-based synthetic agents with empirically grounded behavioral constraints while preserving model transparency and auditability.
\section{Framework}
\subsection{Problem Formulation}
\label{sec:problem_setup}

We formulate \textit{behavioral regularization} as the problem of constraining LLM-generated choice probabilities using empirically estimated behavioral knowledge while preserving the model's ability to reason over rich natural-language policy scenarios. Given a policy scenario represented by an attribute vector $\mathbf{x}$, an unconstrained LLM generates an acceptance probability

\begin{equation}
\pi_{\mathrm{LLM}}(y=1\mid\mathbf{x}),
\end{equation}

where $y\in\{0,1\}$ denotes whether the synthetic respondent accepts the proposed intervention. Although the LLM can flexibly interpret complex policy narratives, the generated probabilities are not guaranteed to satisfy empirically established behavioral relationships.

To provide empirical behavioral guidance, we estimate an empirical choice frontier from discrete choice experiment (DCE) responses using a mixed-logit random utility model~\citep{McFadden1974b,BenAkiva1985}. The resulting frontier,

\begin{equation}
P_{\mathrm{emp}}(y=1\mid\mathbf{x}),
\end{equation}

maps policy-relevant attribute configurations to acceptance probabilities consistent with observed stated preferences. Throughout this paper, we treat this frontier as an empirically estimated behavioral benchmark rather than assumption-free ground truth.

The empirical frontier defines directional behavioral relationships that should be preserved by a behaviorally reliable synthetic agent. For example, for a disutility attribute $x_k$ (e.g., waiting time), the estimated frontier satisfies

\begin{equation}
x_k^{(a)}<x_k^{(b)}
\Longrightarrow
P_{\mathrm{emp}}(y=1\mid x_k^{(a)},\mathbf{x}_{-k})
\ge
P_{\mathrm{emp}}(y=1\mid x_k^{(b)},\mathbf{x}_{-k}),
\label{eq:monotonicity}
\end{equation}

while the inequality is reversed for utility-enhancing attributes such as vaccine efficacy or financial incentives.

We define a \textit{behavioral inconsistency} as a violation of these empirical directional relationships by the LLM-generated probabilities. For example, if two policy scenarios differ only in waiting time and the LLM assigns a higher acceptance probability to the longer-wait scenario, the generated decision violates the empirical behavioral ordering estimated from the DCE.

The objective of Empirical-Frontier Regularization is therefore to construct a behaviorally regularized policy

\[
\pi_{\mathrm{EFR}}(y=1\mid\mathbf{x}),
\]

that combines the representational flexibility of LLMs with empirically estimated behavioral knowledge. Specifically, the framework is designed to achieve three objectives:

\begin{itemize}
\item[(O1)] Improve agreement with the empirical choice frontier.
\item[(O2)] Reduce violations of empirical behavioral constraints, including directional monotonicity.
\item[(O3)] Enable local inference without transmitting individual-level DCE responses to external model providers.
\end{itemize}

The proposed framework consists of three stages (Table~\ref{tab:exp_setup}): (1) empirical behavioral knowledge extraction from DCE responses, (2) unconstrained LLM policy simulation, and (3) inference-time behavioral regularization.

\subsection{Stage 1: Empirical Behavioral Knowledge Extraction}
\label{sec:stage1}

The first stage extracts empirical behavioral knowledge from a vaccination-timing discrete choice experiment (DCE) conducted among 1,027 adults in Wuhan, China. This knowledge is represented as an empirical choice frontier that serves as the behavioral reference for inference-time regularization. Rather than treating the DCE model as the final prediction system, we use it as an interpretable source of empirically estimated behavioral knowledge.

Let

\[
\mathcal{D}_{\mathrm{DCE}}
=
\{(\mathbf{x}_i,y_i)\}_{i=1}^{N},
\]

denote the observed DCE responses, where $\mathbf{x}_i$ is the attribute vector and $y_i\in\{0,1\}$ indicates vaccine acceptance. The empirical choice frontier is estimated using a mixed-logit model under the random utility framework~\citep{McFadden1974b,BenAkiva1985}:

\[
U(\mathbf{x},\boldsymbol{\beta})
=
\beta_0
+
\beta_{\mathrm{wait}}\mathrm{WaitTime}
+
\beta_{\mathrm{eff}}\mathrm{Efficacy}
+
\beta_{\mathrm{side}}\mathrm{SideEffects}
+
\beta_{\mathrm{cash}}\mathrm{Cash},
\]

with choice probability

\[
P(y=1\mid\mathbf{x})
=
\int
\sigma\!\left(U(\mathbf{x},\boldsymbol{\beta})\right)
f(\boldsymbol{\beta}\mid\boldsymbol{\theta})
\,d\boldsymbol{\beta},
\]

where $\sigma(\cdot)$ is the logistic function. Estimated coefficients are reported in Appendix~\ref{app:dce_attributes}.

\paragraph{Construction of the Empirical Behavioral Reference.}

Inference-time regularization requires a deterministic behavioral reference for every policy scenario. For the main experiments, we therefore construct a $4\times4\times4=64$-state evaluation grid by varying waiting time, vaccine efficacy, and side-effect risk across four levels each while fixing the cash incentive at its reference level. The behavioral reference is computed using the posterior mean of the estimated random coefficients,

\[
U_{\mathrm{ref}}(\mathbf{x})
=
\bar{\beta}_{\mathrm{wait}}\mathrm{WaitTime}
+
\bar{\beta}_{\mathrm{eff}}\mathrm{Efficacy}
+
\bar{\beta}_{\mathrm{side}}\mathrm{SideEffects},
\]

and

\[
P_{\mathrm{emp}}(y=1\mid\mathbf{x})
=
\sigma\!\left(U_{\mathrm{ref}}(\mathbf{x})\right).
\]

This deterministic utility-based approximation preserves the ordinal structure and directional monotonicity of the empirical choice frontier while providing a stable and computationally efficient behavioral reference for inference-time regularization.

Throughout this study, the empirical choice frontier is interpreted as an empirically estimated behavioral reference rather than assumption-free ground truth. Accordingly, all behavioral alignment results are evaluated relative to this empirical benchmark.

\subsection{Stage 2: Inference-Time Behavioral Regularization}
\label{sec:stage2}

Given the empirical behavioral reference extracted in Stage~1, the second stage regularizes LLM-generated choice probabilities at inference time. Rather than retraining or fine-tuning the underlying model, the proposed framework operates entirely after generation, combining the LLM's prediction with empirically estimated behavioral knowledge.

For each policy scenario $\mathbf{x}_g$, the LLM is prompted to simulate a respondent and report an acceptance probability. Let

\[
\pi_{\mathrm{LLM}}^{(r)}(y=1\mid\mathbf{x}_g)
\]

denote the acceptance probability extracted from generation $r$. To reduce stochastic variation, the probabilities are averaged over $R$ independent generations,

\[
\bar{\pi}_{\mathrm{LLM}}(y=1\mid\mathbf{x}_g)
=
\frac{1}{R}
\sum_{r=1}^{R}
\pi_{\mathrm{LLM}}^{(r)}(y=1\mid\mathbf{x}_g).
\]

The regularized probability is then defined as

\[
\pi_{\mathrm{EFR}}(y=1\mid\mathbf{x}_g)
=
\lambda
\bar{\pi}_{\mathrm{LLM}}(y=1\mid\mathbf{x}_g)
+
(1-\lambda)
P_{\mathrm{emp}}(y=1\mid\mathbf{x}_g),
\label{eq:anchored_policy}
\]

where $P_{\mathrm{emp}}$ is the empirical behavioral reference estimated in Stage~1 and $\lambda\in[0,1]$ controls the relative contribution of the LLM and the empirical behavioral knowledge. The limiting cases have intuitive interpretations: $\lambda=1$ corresponds to the unconstrained LLM, whereas $\lambda=0$ reduces the framework to the empirical behavioral model,

\[
\pi_{\mathrm{DCE}}(y=1\mid\mathbf{x}_g)
=
P_{\mathrm{emp}}(y=1\mid\mathbf{x}_g).
\]

Unless otherwise specified, we use $\lambda=0.25$, assigning 25\% weight to the LLM prediction and 75\% weight to the empirical behavioral reference. This conservative specification preserves the flexibility of the LLM while allowing the empirical frontier to dominate the final behavioral prediction. Sensitivity analyses over the full range of $\lambda$ are reported in Appendix~\ref{app:lambda}.

Unlike parameter updating or reinforcement-learning-based alignment, the proposed regularization operates entirely at inference time and therefore leaves the underlying foundation model unchanged. Because the empirical behavioral reference, the raw LLM probability, the anchoring weight, and the final regularized probability are all explicitly observable, the framework remains transparent, auditable, and model-agnostic. It can therefore be applied to either locally deployed or cloud-hosted LLMs without modifying model parameters. The conceptual relationships among Pure-DCE, Uniform Shrinkage, EFR, and unconstrained LLM inference are summarized in Appendix~\ref{app:architecture_comparison}.

\subsection{Stage 3: Policy Decision Support}
\label{sec:stage3}

The final stage applies the behaviorally regularized probability surface to policy evaluation. Once LLM-generated probabilities have been regularized, the resulting response surface can be queried to compare hypothetical policy scenarios, estimate changes in expected acceptance, and prioritize competing interventions.

Let $\mathbf{x}=(x_k,\mathbf{x}_{-k})$ denote a baseline policy scenario, where $x_k$ is the target policy attribute and $\mathbf{x}_{-k}$ represents the remaining attributes. A hypothetical policy intervention changes the target attribute from $x_k$ to $x_k'$ while holding all remaining attributes fixed,

\[
\mathbf{x}'=(x_k',\mathbf{x}_{-k}).
\]

The corresponding policy response is defined as

\[
\Delta_k^{\mathrm{EFR}}(\mathbf{x})
=
\pi_{\mathrm{EFR}}(y=1\mid\mathbf{x}')
-
\pi_{\mathrm{EFR}}(y=1\mid\mathbf{x}),
\]

where $\pi_{\mathrm{EFR}}$ is the behaviorally regularized probability defined in Eq.~\eqref{eq:anchored_policy}. Averaging over the set of comparable policy scenarios $\mathcal{G}_k$ gives the mean policy response,

\[
\widehat{\Delta}_k
=
\frac{1}{|\mathcal{G}_k|}
\sum_{\mathbf{x}_g\in\mathcal{G}_k}
\Delta_k^{\mathrm{EFR}}(\mathbf{x}_g).
\]

For the main experiments, policy responses are evaluated on the same $4\times4\times4=64$ evaluation grid used throughout the study, constructed by varying waiting time, vaccine efficacy, and side-effect risk across four levels each while fixing the cash incentive at its reference level.

Beyond estimating changes in expected acceptance, the framework also supports policy-ranking analysis by comparing the ordering of candidate interventions induced by $\pi_{\mathrm{EFR}}$ with that implied by the empirical behavioral reference. Such rankings are particularly relevant for healthcare decision support, where decision makers often seek to prioritize competing interventions rather than evaluate a single policy option in isolation. Policy-ranking performance is evaluated in the experimental section.

Because policy evaluation operates only on the empirical behavioral reference and stored LLM-generated probabilities, the framework can be deployed locally without transmitting individual-level DCE responses to external model providers, supporting transparent and privacy-preserving healthcare decision support. The resulting policy-response quantities are descriptive summaries of the regularized probability surface and should not be interpreted as causal treatment effects.

The output of Stage~3 consists of policy-response summaries and intervention rankings that support downstream healthcare decision making.



\clearpage

\clearpage
\section{Experimental Evaluation}

\subsection{Experimental Setup}
\label{sec:setup}

The overall experimental design is summarized in Table~\ref{tab:exp_setup}. Empirical behavioral knowledge was extracted from a vaccination-timing discrete choice experiment (DCE) conducted among 1,027 adults in Wuhan, China. Behavioral alignment was evaluated on a fixed 64-state policy grid using three LLMs and three prediction settings. Held-out predictive performance was assessed on an independent respondent-level test set.

\begin{table}[t]
\centering
\caption{Summary of the experimental design.}
\label{tab:exp_setup}
\begin{tabular}{ll}
\toprule
Component & Specification\\
\midrule
DCE respondents & 1,027 adults (Wuhan, China)\\
Training / held-out split & 822 / 205 respondents (80\% / 20\%)\\
Empirical behavioral model & Mixed-logit\\
Primary LLM & Qwen2.5-72B-Instruct\\
Robustness LLMs & DeepSeek V4 Pro; MiroThinker-1.7-mini\\
Evaluation grid & $4\times4\times4=64$ policy states\\
Waiting time & 0, 2, 4, 6 months\\
Vaccine efficacy & 30\%, 50\%, 70\%, 90\%\\
Side-effect risk & 0, 1, 2, 3 symptoms\\
Cash incentive & Fixed at reference level\\
LLM generations & $R=10$ per policy state\\
Main regularization weight & $\lambda=0.25$\\
Evaluation metrics & MSE, MAE, MVR, Spearman's $\rho$, Log-loss, Brier score\\
\bottomrule
\end{tabular}
\end{table}

Detailed implementation specifications, including prompt templates, parsing procedures, and computational cost, are reported in Appendix~\ref{app:implementation}.

\clearpage
\subsection{Behavioral Consistency with the Empirical Frontier}
\label{sec:behavioral_consistency}

\begin{table*}[t]
\centering
\caption{Behavioral consistency with the empirical behavioral reference before and after Empirical-Frontier Regularization (EFR). Lower values of mean squared error (MSE), mean absolute error (MAE), and monotonicity violation rate (MVR) indicate better behavioral consistency. MVR is defined as the proportion of policy-state pairs violating the DCE-implied monotonicity constraint. Higher Spearman's rank correlation ($\rho$) indicates better preservation of the empirical preference ordering.}
\label{tab:behavioral_alignment}

\begin{tabular}{llcccc}
\toprule
Model & Method & MSE $\downarrow$ & MAE $\downarrow$ & MVR (\%) $\downarrow$ & Spearman's $\rho$ $\uparrow$\\
\midrule

Qwen2.5-72B-Instruct
& Raw LLM
& 0.0482
& 0.1878
& 62.5
& -0.024 \\

& EFR
& 0.0030
& 0.0469
& 4.2
& 0.952 \\

\addlinespace

DeepSeek V4 Pro
& Raw LLM
& 0.0395
& 0.1704
& 39.6
& 0.380 \\

& EFR
& 0.0025
& 0.0428
& 0.0
& 0.976 \\

\addlinespace

MiroThinker-1.7-mini
& Raw LLM
& 0.0506
& 0.1901
& 43.8
& 0.219 \\

& EFR
& 0.0032
& 0.0476
& 0.0
& 0.970 \\

\bottomrule
\end{tabular}
\end{table*}

Table~\ref{tab:behavioral_alignment} summarizes behavioral consistency before and after Empirical-Frontier Regularization (EFR). Across all evaluated language models, inference-time regularization substantially improved agreement with the empirical behavioral reference. The pilot experiment used to develop the evaluation framework is reported in Appendix~\ref{app:pilot}.

For the primary model, Qwen2.5-72B-Instruct, EFR reduced the frontier-deviation mean squared error (MSE) from 0.0482 to 0.0030 and the mean absolute error (MAE) from 0.1878 to 0.0469. The largest improvement was observed in structural behavioral consistency: the waiting-time monotonicity violation rate (MVR) decreased from 62.5\% to 4.2\%, where a violation indicates that the model assigns a higher acceptance probability to a longer waiting time than to a shorter one. Spearman's rank correlation simultaneously increased from $-0.024$ to 0.952, indicating substantially improved preservation of the empirical preference ordering.

The same qualitative pattern was observed for the two additional language models. For DeepSeek V4 Pro, EFR reduced MSE from 0.0395 to 0.0025, eliminated all observed waiting-time monotonicity violations (MVR = 0\%), and increased Spearman's $\rho$ from 0.380 to 0.976. MiroThinker-1.7-mini exhibited similar improvements, with MSE decreasing from 0.0506 to 0.0032, MVR decreasing from 43.8\% to 0\%, and Spearman's $\rho$ increasing from 0.219 to 0.970.

Across all evaluated models, EFR consistently improved both numerical agreement with the empirical behavioral reference (MSE and MAE) and structural behavioral consistency (MVR and Spearman's $\rho$).

\begin{figure}[t]
\centering
\includegraphics[width=0.85\textwidth]{fig2_wait_time_monotonicity_repair.png}
\caption{
Representative illustration of Empirical-Frontier Regularization (EFR) for the waiting-time attribute. The dashed curve represents the empirical behavioral reference estimated from the DCE. Red points denote probabilities generated by the unconstrained Qwen2.5-72B model, and blue points denote the EFR-regularized probabilities ($\lambda=0.25$). The unconstrained LLM assigns higher acceptance probabilities to several longer waiting-time states, violating the empirical monotonicity constraint. EFR substantially restores the expected monotonic ordering by shifting the predicted probabilities toward the empirical behavioral reference while preserving the overall response pattern.
}
\label{fig:frontier_anchor}
\end{figure}

Figure~\ref{fig:frontier_anchor} provides a representative visualization of the regularization process. Consistent with the quantitative results in Table~\ref{tab:behavioral_alignment}, EFR substantially reduces frontier-relative behavioral inconsistencies without modifying the underlying language model. Rather than retraining or fine-tuning the LLM, EFR operates entirely at inference time by combining the unconstrained prediction with the empirical behavioral reference. The resulting probability surface more closely follows the DCE-implied preference ordering while preserving the model's ability to interpret the policy scenario.

\clearpage
\subsection{Held-Out Predictive Validation}
\label{sec:heldout_validation}

The frontier-alignment results in Section~\ref{sec:behavioral_consistency} are benchmark-relative because the anchored probabilities are constructed partly from the empirical choice frontier. We therefore evaluated predictive performance on an independent respondent-level held-out test set using the 80/20 split described in Section~\ref{sec:setup}. To prevent information leakage, the reference model was estimated exclusively from the training respondents. Because LLM probabilities were available only for the predefined 64-state evaluation grid, comparisons involving the LLM and BDT were restricted to the 809 held-out observations whose attribute profiles matched the evaluation grid exactly. Covariate balance diagnostics (Appendix~E.10) indicate that this matched subset systematically differs from the full held-out sample on all observed attributes (SMD~$>0.6$ for waiting time and side effects), consistent with the grid's restriction to specific integer attribute levels. The matched subset should therefore be interpreted as a non-representative diagnostic sample rather than a generalizable test set.

\begin{table}[!t]
\centering
\caption{Held-out predictive performance on DCE choice data.}
\label{tab:heldout}
\small
\resizebox{\textwidth}{!}{%
\begin{tabular}{lccccc}
\toprule
Method & $N$ & Log loss $\downarrow$ & Brier $\downarrow$ & Calibration intercept & Calibration slope \\
\midrule
\multicolumn{6}{l}{\textit{Panel A: Full held-out test sample}}\\
Pure-DCE & 3,690 & 0.604 & 0.208 [0.201, 0.214] & $-0.125$ & 0.818 \\
\midrule
\multicolumn{6}{l}{\textit{Panel B: LLM-matched held-out subset}}\\
Pure-DCE & 809 & 0.688 [0.676, 0.700] & 0.247 [0.241, 0.252] & $-0.201$ & 0.680 \\
Unconstrained LLM & 809 & 0.815 [0.794, 0.836] & 0.308 [0.298, 0.318] & $-0.554$ & $-2.218$ \\
BDT ($\lambda=0.25$) & 809 & 0.703 [0.692, 0.714] & 0.255 [0.249, 0.260] & 0.810 [0.678, 0.953] & 1.523 [1.347, 1.720] \\
\bottomrule
\end{tabular}
}
\vspace{0.4em}

\begin{minipage}{\textwidth}
\small
\textit{Note.} Panel A evaluates the train-only conditional logit benchmark on the full held-out test sample. Panel B is restricted to held-out observations represented in the 64-state evaluation grid. BDT probabilities were computed as $\pi_{\mathrm{BDT}}=0.25\bar{\pi}_{\mathrm{LLM}}+0.75P_{\mathrm{static}}$. Bracketed values denote 95\% respondent-level cluster bootstrap confidence intervals (2,000 replications).
\end{minipage}
\end{table}

Table~\ref{tab:heldout} shows that the train-only Pure-DCE model provides the strongest predictive benchmark within the structured DCE attribute space, achieving a log loss of 0.604 and a Brier score of 0.208 on the full held-out test sample. This result is expected because the evaluation remains within the original DCE design and therefore favors the model directly estimated from the observed stated-preference data.

On the LLM-matched held-out subset, BDT substantially improves predictive performance relative to the unconstrained LLM. Log loss decreases from 0.815 to 0.703, and the Brier score decreases from 0.308 to 0.255. The respondent-level bootstrap intervals for both metrics do not overlap between the unconstrained LLM and BDT, indicating a statistically meaningful improvement in predictive accuracy.

Pure-DCE remains slightly more accurate than BDT on the matched subset (log loss: 0.688 vs.\ 0.703; Brier score: 0.247 vs.\ 0.255), indicating that the structured choice model remains the predictive ceiling within the original DCE attribute space. The performance gap, however, is modest, suggesting that BDT preserves most of the predictive accuracy of the empirical benchmark while retaining the flexibility of an LLM for richer policy scenarios.

Calibration diagnostics reveal a complementary pattern. Relative to the unconstrained LLM (calibration slope $=-2.218$), BDT restores the correct calibration direction (slope $=1.523$) and substantially improves probability prediction. However, the calibration intercept remains positive and the slope exceeds one, indicating residual over-confidence introduced by the retained LLM component. This behavior is expected because the anchoring rule intentionally preserves a non-zero contribution from the unconstrained LLM ($\lambda=0.25$). Future work may explore adaptive or attribute-specific anchoring weights to further improve probabilistic calibration.

Overall, the held-out validation demonstrates that BDT substantially narrows the predictive gap between unconstrained LLMs and the empirical DCE benchmark while preserving the behavioral consistency established in Section~\ref{sec:behavioral_consistency}.

\clearpage

\clearpage
\subsection{Out-of-Design Policy Scenario Evaluation}
\label{sec:out_of_design}

The evaluations in Sections~\ref{sec:behavioral_consistency} and \ref{sec:heldout_validation} were conducted within the structured DCE attribute space. We next examine whether empirical-frontier regularization preserves behavioral consistency when LLMs are applied to richer policy narratives that extend beyond the original DCE design. Unlike the previous experiments, these scenarios require the LLM to interpret unstructured textual descriptions while maintaining the waiting-time monotonicity implied by the empirical behavioral reference.

Twenty paired policy scenarios were constructed and divided into two evaluation sets. Panel A (six paired scenarios) introduced common access-related descriptions not explicitly represented in the DCE, including walk-in vaccination, online registration, and administrative paperwork. Panel B (fourteen paired scenarios) introduced richer policy narratives containing fictional incentives (e.g., VIP access and cash rewards) together with multi-attribute descriptions designed to challenge behavioral consistency beyond the original DCE design. Across all scenarios, waiting time remained the target attribute used to evaluate monotonicity, whereas the additional textual descriptions were interpreted only by the LLM component.

\begin{table}[t]
\centering
\caption{Behavioral consistency under out-of-design policy scenarios.}
\label{tab:out_of_design}
\small
\setlength{\tabcolsep}{4pt}
\begin{tabular}{@{}llccc@{}}
\toprule
Model & Method & MVR ($\downarrow$) & MAD ($\downarrow$) & Spearman $\rho$ ($\uparrow$)\\
\midrule

\multicolumn{5}{l}{\textit{Panel A: Core scenarios (6 paired policy descriptions)}}\\

Qwen-max
& Unconstrained LLM
& 0.333 (2/6)
& 0.290
& 0.12 \\

Qwen-max
& BDT
& 0.000
& 0.072
& 0.87 \\

DeepSeek V3
& Unconstrained LLM
& 0.667 (4/6)
& 0.257
& $-0.07$ \\

DeepSeek V3
& BDT
& 0.000
& 0.064
& 0.87 \\

\midrule

\multicolumn{5}{l}{\textit{Panel B: Extended scenarios (14 paired policy descriptions)}}\\

Qwen-max
& Unconstrained LLM
& 0.286 (4/14)
& 0.125
& $-0.01$ \\

Qwen-max
& BDT
& 0.000
& 0.031
& 0.93 \\

DeepSeek V3
& Unconstrained LLM
& 0.357 (5/14)
& 0.209
& $-0.06$ \\

DeepSeek V3
& BDT
& 0.000
& 0.052
& 0.88 \\

\bottomrule
\end{tabular}

\vspace{0.4em}

\begin{minipage}{0.95\textwidth}
\small
\textit{Note.}
Twenty paired policy scenarios were constructed. Panel A evaluates basic access-related policy descriptions, whereas Panel B introduces richer textual narratives containing fictional incentives and multi-attribute descriptions. MVR denotes the monotonicity violation rate across paired scenarios, MAD denotes the mean absolute deviation from the empirical behavioral reference, and Spearman's $\rho$ measures rank agreement with the empirical behavioral reference. Because Panel B contains fourteen paired scenarios, these statistics should be interpreted as descriptive robustness measures rather than population-level estimates. This experiment evaluates behavioral consistency under out-of-design policy narratives rather than predictive accuracy against observed outcomes.
\end{minipage}
\end{table}

Table~\ref{tab:out_of_design} summarizes the out-of-design evaluation. In the core scenarios (Panel A), the unconstrained LLMs violated the expected waiting-time ordering in 2 of 6 paired scenarios for Qwen-max and 4 of 6 paired scenarios for DeepSeek V3. After empirical-frontier regularization, no monotonicity violations were observed for either model, and rank agreement with the empirical behavioral reference increased substantially.

The differences became more pronounced in the extended scenarios (Panel B). Rich textual narratives containing fictional incentives and complex policy descriptions reduced the rank agreement of the unconstrained LLMs to approximately zero for both models. In contrast, BDT maintained perfect monotonicity (0/14 violations) while restoring rank agreement to $\rho=0.93$ for Qwen-max and $\rho=0.88$ for DeepSeek V3. These results indicate that empirical-frontier regularization continues to preserve the waiting-time monotonicity implied by the empirical behavioral reference even when language models are required to interpret policy narratives beyond the original DCE design.

This experiment evaluates the behavioral consistency of the final anchored predictions rather than isolating the respective contributions of language understanding and empirical regularization. Additional robustness analyses examining the influence of DCE sample size on empirical-frontier estimation are reported in Appendix~\ref{app:dce_sample_sensitivity}.
\clearpage
\subsection{Policy Ranking Consistency}
\label{sec:policy_ranking}

While the preceding analyses evaluated state-level behavioral consistency, many policy applications depend primarily on whether competing interventions are ranked in the correct order. We therefore evaluated Policy Ranking Consistency (PRC), defined in Section~3.4, using the Pure-DCE model as the empirical behavioral reference. Six candidate policy interventions generated 15 pairwise comparisons, and PRC was computed as the proportion of policy pairs whose ordering agreed with the Pure-DCE ranking.

\begin{table}[t]
\centering
\caption{Policy Ranking Consistency (PRC) relative to the Pure-DCE reference.}
\label{tab:prc}
\small
\setlength{\tabcolsep}{5pt}
\begin{tabular}{l l c c c}
\toprule
Model & Method & Matching pairs & PRC & 95\% CI \\
\midrule
Qwen2.5-72B & Unconstrained LLM & 5/15 & 0.333 & [0.118, 0.616] \\
Qwen2.5-72B & EFR & 13/15 & 0.867 & [0.595, 0.983] \\
\midrule
DeepSeek V4 Pro & Unconstrained LLM & 2/15 & 0.133 & [0.017, 0.405] \\
DeepSeek V4 Pro & EFR & 9/15 & 0.600 & [0.323, 0.837] \\
\midrule
MiroThinker-1.7-mini & Unconstrained LLM & 8/15 & 0.533 & [0.266, 0.787] \\
MiroThinker-1.7-mini & EFR & 12/15 & 0.800 & [0.519, 0.957] \\
\bottomrule
\end{tabular}

\vspace{0.4em}
\begin{minipage}{0.95\textwidth}
\small
\textit{Note.} PRC measures the proportion of intervention pairs ranked in the same order as the Pure-DCE reference. The policy set contains six interventions, yielding 15 unordered policy pairs. Brackets report exact binomial 95\% confidence intervals. One-sided exact binomial tests were performed against the random-ranking baseline (PRC = 0.5). Benjamini--Hochberg correction was applied across the three EFR model comparisons (Qwen2.5-72B, DeepSeek V4 Pro, and MiroThinker-1.7-mini); Qwen2.5-72B and MiroThinker-1.7-mini remained significant after correction.
\end{minipage}
\end{table}

Table~\ref{tab:prc} shows that EFR consistently improved policy-ranking agreement across all three language models. For Qwen2.5-72B, PRC increased from 0.333 to 0.867, correctly ranking 13 of the 15 intervention pairs. Notably, the unconstrained Qwen2.5-72B model produced a PRC below the random-ranking baseline (0.333 < 0.5), indicating systematic rather than random disagreement with the empirical policy ordering. This pattern is consistent with the state-level behavioral violations observed in Section~4.2: systematic reversals of monotonicity naturally propagate to incorrect rankings of policy interventions. EFR largely corrected these systematic ranking errors. MiroThinker-1.7-mini similarly improved from 0.533 to 0.800, whereas DeepSeek V4 Pro improved from 0.133 to 0.600.

The strongest evidence was observed for Qwen2.5-72B and MiroThinker-1.7-mini. For both models, the confidence intervals remained above the random-ranking baseline, indicating statistically significant agreement with the Pure-DCE policy ordering after multiple-testing correction. DeepSeek V4 Pro also exhibited a substantial improvement in point estimates, although its confidence interval remained wider and included the null value, indicating greater uncertainty in the estimated ranking consistency.

These findings complement the behavioral alignment results reported in Section~4.2. Beyond improving agreement with the empirical behavioral reference at the individual state level, EFR also produces policy rankings that more closely match those implied by the DCE model. This property is particularly relevant for decision support because many policy analyses rely primarily on the relative ordering of competing interventions rather than on precise estimates of absolute uptake probabilities.

The interpretation of PRC remains benchmark-relative. The reference ranking is derived from the DCE-estimated empirical behavioral reference rather than observed policy outcomes. Consequently, PRC evaluates the extent to which LLM-based policy simulations preserve the intervention ordering implied by the empirical choice model, while external validation is still required to establish real-world policy predictive validity.
The policy simulation framework, uncertainty intervals, and additional validation details are provided in Appendices~\ref{app:pure_dce_prc} and \ref{app:bootstrap}.


\section{Discussion}
\subsection{Principal Findings}

This study demonstrates that inference-time empirical regularization can substantially improve the behavioral consistency of LLM-based synthetic agents without retraining the underlying language model. Across three open-source LLMs, Empirical-Frontier Regularization (EFR) consistently aligned generated choice probabilities with an empirically estimated DCE behavioral reference, reducing waiting-time monotonicity violations from 62.5\% to 4.2\% for the primary Qwen2.5-72B model while improving rank correlation from $-0.024$ to 0.952.

Three principal findings emerge. First, unconstrained LLMs frequently produced probabilities that violated empirically estimated preference structure despite generating linguistically coherent explanations. These systematic behavioral inconsistencies propagated from individual choice states to policy-level intervention rankings. By introducing an explicit empirical behavioral reference during inference, EFR substantially reduced these systematic deviations while preserving the LLM's ability to interpret richer policy narratives.

Second, empirical regularization improved predictive performance relative to unconstrained LLMs while maintaining most of the predictive accuracy of the structured DCE benchmark. Held-out validation showed lower log loss and Brier score than the unconstrained models, although the Pure-DCE model remained the strongest predictor within the original DCE attribute space. At the policy level, EFR also generated intervention rankings that more closely matched those implied by the empirical choice model, suggesting that improved state-level behavioral consistency translated into more reliable policy comparisons.

Taken together, these findings suggest that the principal limitation of unconstrained LLM synthetic agents is not their capacity to understand complex policy narratives, but the absence of explicit empirical constraints during inference. Rather than replacing either LLMs or discrete choice models, EFR combines their complementary strengths by coupling flexible language understanding with empirically grounded behavioral regularization.

\subsection{Methodological Contributions}

The primary methodological contribution of this study is the introduction of inference-time empirical regularization as a transparent mechanism for aligning LLM-generated decisions with empirically estimated behavioral structure. Existing approaches to LLM alignment, such as instruction tuning, reinforcement learning from human feedback (RLHF), and prompt engineering, primarily optimize linguistic quality, helpfulness, or pairwise preference judgments. In contrast, EFR targets a different objective: preserving structural consistency with an empirically estimated choice model. Rather than modifying model parameters or retraining the underlying LLM, EFR regularizes only the generated probability surface through a post-generation convex anchoring rule. More broadly, this design reflects a principle of \emph{theory-informed machine learning}, in which established empirical models provide transparent structural priors for generative AI systems.

Importantly, EFR should not be viewed as an alternative to either discrete choice models or existing LLM alignment methods. Instead, it combines the complementary strengths of both approaches. The empirical choice model provides an explicit behavioral reference within experimentally observed attribute dimensions, whereas the LLM contributes flexible reasoning over richer textual scenarios that extend beyond the original DCE design. By coupling these components at inference time, EFR preserves the interpretability and auditability of the empirical model while retaining the generative flexibility of the language model.

Viewed from this perspective, EFR represents a practical example of theory-informed machine learning. Rather than relying exclusively on end-to-end data-driven optimization, the framework incorporates an empirically estimated behavioral model directly into the inference process. The empirical behavioral reference functions as an explicit structural prior that constrains generated decisions on dimensions supported by stated-preference evidence while allowing the LLM to interpret unstructured contextual information. Unlike conventional alignment methods, these behavioral constraints remain transparent, independently estimable, and directly auditable throughout the inference process.

The broader significance of EFR extends beyond the specific vaccination case study presented here. In principle, similar architectures could combine foundation models with clinical risk prediction models, survival models, causal effect estimators, or other interpretable statistical models in settings where a well-estimated empirical model is available. Although these applications remain to be demonstrated empirically, the central idea is general: established domain models can provide explicit inference-time constraints that complement the flexibility of foundation models. Table~\ref{tab:architecture_comparison} in Appendix~A provides a conceptual comparison of the four inference-time architectures evaluated in this study, illustrating their respective positions along the dimensions of LLM heterogeneity, transparency, and behavioral consistency. Viewed in this way, EFR represents not merely a method for improving synthetic-agent behavior, but a general framework for integrating transparent empirical models with generative AI in biomedical informatics.

\subsection{Implications for Biomedical Informatics}

The present work has broader implications for the development of trustworthy AI systems in biomedical informatics. Increasingly, foundation models are being incorporated into clinical decision support, patient communication, and healthcare simulation. However, these applications require more than fluent language generation; they also require generated recommendations and simulated decisions to remain consistent with established empirical evidence. Our results suggest that inference-time empirical regularization provides one practical mechanism for achieving this objective while enabling deployment pathways that avoid the computational and data requirements of model fine-tuning.

More generally, EFR illustrates a hybrid paradigm in which foundation models and conventional statistical models play complementary rather than competing roles. In the proposed framework, the LLM is responsible for interpreting rich textual context and generating flexible responses, whereas the empirical behavioral model provides transparent constraints over clinically or behaviorally meaningful dimensions. This division of responsibilities allows each component to contribute its respective strengths while preserving interpretability, empirical accountability, and flexibility for complex decision-support scenarios.

The framework also aligns with the growing emphasis on trustworthy and auditable AI in healthcare. Because the empirical behavioral reference, anchoring weight, and final predicted probabilities remain explicitly observable, the complete inference process can be independently inspected and reproduced. Such transparency may be particularly valuable in biomedical settings where model outputs require clinical justification, regulatory review, or post hoc auditing before deployment.

Although this study focuses on DCE-estimated behavioral models, the underlying principle is more general. In settings where a well-estimated empirical model is available, inference-time regularization could, in principle, integrate foundation models with clinical risk prediction models, survival models, causal effect estimators, or other interpretable statistical frameworks. We therefore view EFR not only as a framework for improving synthetic-agent behavior, but also as a broader strategy for integrating empirical knowledge with generative AI in biomedical informatics.

\subsection{Limitations and Future Directions}

Several limitations should be considered when interpreting the present findings. First, EFR is inherently constrained by the quality of its empirical behavioral reference. The DCE-estimated frontier represents stated preferences rather than observed behavior and may therefore reflect hypothetical bias, attribute-range limitations, or model misspecification. Consequently, EFR cannot correct biases that are already embedded in the empirical reference.

Second, the current implementation adopts a single global anchoring weight ($\lambda=0.25$) and a population-level behavioral reference. Although this simple specification proved effective across multiple evaluation settings, the held-out calibration analysis suggests that a uniform anchoring weight may not be optimal across the entire probability space. Future work should investigate adaptive anchoring strategies, including attribute-specific or respondent-specific regularization based on posterior preference heterogeneity.

Third, the empirical evaluation was conducted using a single vaccination-timing DCE collected in Wuhan, China. Although the underlying inference-time regularization framework is model-agnostic, the generalizability of the empirical findings to other healthcare decisions, patient populations, and policy contexts remains to be established through external validation. Similarly, while the framework consistently improved behavioral consistency across three open-source LLMs, the magnitude of improvement varied across model families, suggesting that the effectiveness of empirical regularization may depend on characteristics of the underlying language model.

Finally, the present study should be viewed as a proof-of-concept demonstration of empirically constrained LLM simulation rather than a complete, individualized behavioral model. The current framework regularizes population-level response surfaces and does not yet represent respondent-specific preference dynamics or longitudinal behavioral adaptation. Extending inference-time empirical regularization to personalized behavioral models, multiple empirical datasets, and longitudinal decision processes represents an important direction for future research.

Taken together, these limitations suggest that future work should focus on broader external validation, adaptive empirical regularization strategies, and integration with richer longitudinal and individualized behavioral models. Such developments will help determine the extent to which empirically grounded foundation models can support trustworthy decision making across a wider range of biomedical informatics applications.

\subsection{Conclusion}
\label{sec:conclusion}

This study proposed Empirical-Frontier Regularization (EFR), an inference-time framework for aligning LLM-generated behavioral simulations with empirically estimated discrete choice models. Rather than treating large language models as standalone behavioral simulators, EFR incorporates structured behavioral evidence during inference, allowing flexible language understanding to be combined with empirically grounded behavioral constraints.

Across three open-source language models, EFR consistently improved behavioral consistency relative to unconstrained LLMs. The framework reduced frontier-relative inconsistencies, improved held-out probability prediction, and generated policy rankings that more closely matched the empirical behavioral reference while preserving the flexibility of language models for richer textual scenarios. These findings demonstrate that inference-time empirical regularization can substantially improve the consistency of LLM-based synthetic agents without retraining the underlying model.

More broadly, this work demonstrates that hybrid architectures combining generative AI with interpretable empirical models represent a practical and scientifically grounded pathway toward trustworthy AI in biomedical informatics. EFR is not intended to replace either discrete choice models or foundation models; rather, it illustrates how their complementary strengths can be combined through transparent inference-time regularization. We hope this framework provides a foundation for future research on empirically grounded generative AI for biomedical informatics, clinical decision support, and policy simulation.

\clearpage

\bibliographystyle{plainnat}
\bibliography{references}

\clearpage

\appendix

\section{DCE Attribute Levels and Static Frontier Estimates}
\label{app:dce_attributes}

This appendix documents the attribute structure of the Wuhan vaccination-timing discrete choice experiment (DCE) and the static random utility estimates used to construct the empirical choice frontier in the main analysis. The full DCE instrument varied four policy-relevant attributes: waiting time, vaccine efficacy, side-effect risk, and cash incentives. The main 64-state evaluation grid used in the full-scale BDT experiments focuses on three core attributes: waiting time, vaccine efficacy, and side-effect burden. These three dimensions define the policy grid on which the unconstrained LLM and Static-BDT Anchor are evaluated in Table~\ref{tab:behavioral_alignment}.

\subsection{DCE Attribute Levels}

Table~\ref{tab:dce_attributes} reports the attribute levels used in the DCE. Waiting time and side effects are disutility attributes: higher values are expected to lower vaccine acceptance. Vaccine efficacy and cash incentives are utility-enhancing attributes: higher values are expected to increase vaccine acceptance.

\begin{table}[htbp]
\centering
\caption{DCE attribute levels.}
\label{tab:dce_attributes}
\begin{tabular}{lll}
\hline
Attribute & Levels & Expected direction \\
\hline
Wait Time (months) & 0, 1, 2, 3, 6 & Negative \\
Vaccine Efficacy & 50\%, 70\%, 90\%, 95\% & Positive \\
SideEffects & 1\%, 5\%, 10\% & Negative \\
Cash Incentives (RMB) & 0, 200, 500, 800 & Positive \\
\hline
\end{tabular}
\end{table}

\subsection{Static Empirical Choice Frontier}

The empirical frontier used in the main full-scale experiment is constructed using a \textbf{mixed-logit (random parameters logit)} specification to account for unobserved preference heterogeneity across respondents. 

For the construction of the 64-state evaluation grid and the computation of the Static-BDT anchor, the systematic utility for alternative $x$ is strictly specified as:
\[
U(x, \beta) = \beta_{wait} \cdot \text{WaitTime} + \beta_{eff} \cdot \text{Efficacy} + \beta_{side} \cdot \text{SideEffects}.
\]
Crucially, in this grid-level evaluation, the Cash Incentives attribute is held at its reference value (0 RMB), and no opt-out constant (intercept) is included in the utility function, as the grid focuses purely on the trade-offs between the three focal attributes.

Unlike the standard multinomial logit model, the mixed-logit choice probability is given by the integral over the distribution of the random parameters $f(\beta|\theta)$. In our specification, only the coefficient for Wait Time ($\beta_{wait}$) is assumed to follow a continuous distribution (e.g., normal), capturing the heterogeneity in sensitivity to waiting time across the population, while $\beta_{eff}$ and $\beta_{side}$ are treated as fixed parameters. The grid-level benchmark is computed using the mean of the estimated random parameters (reported in Table~\ref{tab:static_frontier_estimates}) as a deterministic approximation:
\[
P_{\mathrm{static}}(y=1|x) \approx \sigma(U(x, \bar{\beta})) = \frac{1}{1 + \exp[-U(x, \bar{\beta})]},
\]
where $\bar{\beta}$ is the vector of mean coefficients. This approximation ensures a consistent, computationally stable, and strictly monotonic behavioral benchmark for the post-generation anchoring rule.

\subsection{Random Utility Estimates}

Table~\ref{tab:static_frontier_estimates} reports the static random utility estimates used to construct the empirical choice frontier. The signs follow the expected behavioral directions: waiting time and side effects reduce acceptance, while vaccine efficacy and cash incentives increase acceptance.

\begin{table}[h]
\centering
\caption{Mixed-logit random utility estimates used to construct the empirical choice frontier.}
\label{tab:static_frontier_estimates}
\small
\setlength{\tabcolsep}{8pt} % 增加列间距，提升可读性
\begin{tabular}{l c c c}
\toprule
\textbf{Attribute} & \textbf{Mean Coeff.} & \textbf{Std. Dev. (if random)} & \textbf{S.E. (Mean)} \\
\midrule
Wait Time & $-0.267^{***}$ & $0.915^{***}$ & $(0.032)$ \\
Vaccine Efficacy & $1.459^{***}$ & --- & $(0.164)$ \\
Side Effects & $-0.102^{***}$ & --- & $(0.037)$ \\
Cash Incentives & $0.003^{***}$ & --- & $(0.0003)$ \\
\bottomrule
\end{tabular}

\vspace{0.4em}
\begin{minipage}{0.95\textwidth}
\small
\textit{Note.} Estimates are from the mixed-logit (random parameters logit) model actually used to compute the 64-state grid probabilities ($P_{\mathrm{static}}$). Mean coefficients represent the expected marginal utility. The ``Std. Dev.'' column reports the standard deviation of the random parameters; only Wait Time was specified as random in the final estimation. Attributes marked with ``---'' are treated as fixed parameters. Standard errors for the mean coefficients are in parentheses. Cash Incentives are included in the full DCE utility estimation but are held at their reference value (0 RMB) in the 64-state grid evaluation. No opt-out constant (intercept) is included in the grid-level utility specification. $^{***}p < 0.001$.
\end{minipage}
\end{table}

All coefficients display the expected sign. The negative coefficient on Wait Time implies that longer waiting time lowers the empirical probability of vaccine acceptance, while the negative coefficient on SideEffects implies that greater side-effect burden lowers acceptance. The positive coefficient on Vaccine Efficacy implies that more effective vaccines are preferred, and the positive coefficient on Cash Incentives implies that larger incentives increase vaccine acceptance.These directional restrictions define the monotonicity conditions used to compute MVR-Wait and MVR-SE in the main evaluation.

\subsection{Connection to the 64-State Evaluation Grid}

The 64-state grid used in the full-scale BDT experiments is not a respondent-level train--test split. Instead, it is a fixed full-factorial policy grid over the three core attributes:

\[
G
=
\{(\text{wait}, \text{eff}, \text{se}) :
\text{wait} \in \{0,2,4,6\},
\text{eff} \in \{0.30,0.50,0.70,0.90\},
\text{se} \in \{0,1,2,3\}
\}.
\]

For each state \(x_g \in G\), the empirical benchmark is

\[
P_{\text{static},g}
=
P_{\text{static}}(y=1 \mid x_g).
\]

The full-scale LLM and Static-BDT results in Table~\ref{tab:behavioral_alignment} are evaluated by comparing agent-implied probabilities against \(P_{\text{static},g}\) over this common 64-state grid. This design ensures that all model comparisons are made on the same policy-relevant attribute space.

\subsection{Comparison of Inference-Time Regularization Architectures}
\label{app:architecture_comparison}

To clarify the design rationale of the proposed framework, Table~\ref{tab:architecture_comparison} summarizes the conceptual differences among the four inference-time architectures evaluated in this study. The comparison emphasizes architectural properties rather than predictive performance.

Pure-DCE represents the empirical behavioral reference estimated directly from the DCE model. Uniform Shrinkage replaces the LLM output with a global constant before convex anchoring, providing a baseline that isolates the effect of parametric shrinkage toward the empirical frontier. EFR extends this formulation by retaining the LLM-generated probability within the same transparent convex anchoring framework. Finally, the unconstrained LLM relies entirely on its internal representations without explicit empirical regularization.

This comparison illustrates the distinct roles of empirical behavioral models and foundation models within hybrid AI systems. The empirical frontier provides the primary source of behavioral discipline, while the LLM supplies the representational capacity required to process richer natural-language policy scenarios. The proposed EFR architecture explicitly separates these two functions through a transparent and auditable inference-time design.

\begin{table}[htbp]
\centering
\caption{Conceptual comparison of inference-time architectures for LLM-based behavioral simulation.}
\label{tab:architecture_comparison}
\small
\setlength{\tabcolsep}{3pt}
\begin{tabular}{%
>{\raggedright\arraybackslash}p{0.29\textwidth}
>{\centering\arraybackslash}p{0.10\textwidth}
>{\centering\arraybackslash}p{0.13\textwidth}
>{\centering\arraybackslash}p{0.10\textwidth}
>{\centering\arraybackslash}p{0.07\textwidth}
>{\centering\arraybackslash}p{0.09\textwidth}
>{\centering\arraybackslash}p{0.15\textwidth}}
\toprule
Method &
Retains LLM variation &
Transparent \& auditable &
Inference-time &
MVR-Wait $\downarrow$ &
$\rho$ $\uparrow$ &
Supports rich textual scenarios \\
\midrule
Pure-DCE\newline
{\footnotesize ($P_{\mathrm{static}}$)}
& $\times$ & $\checkmark$ & $\checkmark$ & 0.000 & 1.000 & $\times$ \\

Uniform Shrinkage\newline
{\footnotesize ($\lambda\bar{p}+(1-\lambda)P_{\mathrm{static}}$)}
& $\times$ & $\checkmark$ & $\checkmark$ & 0.000 & 1.000 & $\times$ \\

EFR (this work)\newline
{\footnotesize ($\lambda\hat{p}_{\mathrm{LLM}}+(1-\lambda)P_{\mathrm{static}}$)}
& $\checkmark$ & $\checkmark$ & $\checkmark$ & 0.042 & 0.952 & $\checkmark$ \\

Unconstrained LLM
& $\checkmark$ & $\times$ & $\checkmark$ & 0.625 & $-0.024$ & $\checkmark$ \\
\bottomrule
\end{tabular}

\vspace{0.4em}
\begin{minipage}{0.96\textwidth}
\small
\textit{Note.} Metrics are shown for the primary Qwen2.5-72B model under the main specification ($\lambda=0.25$). Uniform Shrinkage replaces the LLM probability with the global mean acceptance probability, isolating the effect of parametric shrinkage toward the empirical frontier. "Supports rich textual scenarios" indicates whether the architecture can incorporate information from unstructured natural-language policy descriptions during inference; it should not be interpreted as evidence of superior predictive performance in out-of-design settings.
\end{minipage}

\end{table}


\section{Pilot Experiment Details}
\label{app:pilot}

This appendix documents the pilot experiment reported in the Pilot rows of Table~\ref{tab:behavioral_alignment}. The pilot served as a low-resource mechanism check rather than a directly comparable full-scale evaluation. It tested whether explicit monotonicity information could reduce frontier-relative behavioral inconsistencies in a minimal setting before scaling to the full Static-BDT Anchor implementation.

\subsection{Pilot Design}

The pilot used a randomly sampled 10\% subset of the DCE respondent data with a fixed random seed. The base model was Qwen2.5-1.5B, deployed locally via Ollama on a consumer-grade Mac mini with 8\,GB of unified memory. The pilot evaluated two conditions:

\begin{enumerate}
    \item \textbf{Unconstrained LLM}: the model was prompted as a patient considering a vaccine, without any explicit behavioral rule.
    \item \textbf{Rule-injection BDT}: the model was prompted with a simple monotonicity rule derived from the DCE logic: longer waiting time should reduce willingness to accept vaccination.
\end{enumerate}

Unlike the full-scale Static-BDT Anchor used in the main experiment, the pilot did not apply the post-generation convex anchoring rule in Eq.~\eqref{eq:anchored_policy}. Instead, it used prompt-level rule injection. The pilot rows in Table~\ref{tab:behavioral_alignment} should therefore be interpreted as low-resource mechanism evidence that explicit monotonicity information can reduce frontier-relative behavioral inconsistencies, not as a direct estimate of the full frontier-anchored calibration mechanism.

\subsection{Pilot Prompt Templates}

For reproducibility, the pilot used two prompt templates (unconstrained baseline and rule-injection BDT), each presenting a vaccination scenario with varied waiting time while holding cost and side-effect risk fixed. The unconstrained prompt asked the model to decide without additional constraints. The rule-injection variant added an explicit monotonicity rule ("Longer waiting time decreases willingness to accept") before the same scenario. Complete prompt templates are provided in the replication materials.

\subsection{Pilot MVR Calculation}

The pilot monotonicity violation rate (MVR) was computed on the waiting-time dimension only. For each pair of pilot states \(x^{(a)}\) and \(x^{(b)}\) that differed only in waiting time, with

\[
\text{WaitTime}^{(a)} < \text{WaitTime}^{(b)},
\]

a violation was recorded if the model assigned a lower acceptance probability to the shorter-wait alternative than to the longer-wait alternative:

\[
\bar{\pi}(y=1 \mid x^{(a)})
<
\bar{\pi}(y=1 \mid x^{(b)}).
\]

The pilot MVR is therefore

\[
\text{MVR}
=
\frac{
\# \{(x^{(a)},x^{(b)}): \text{WaitTime}^{(a)} < \text{WaitTime}^{(b)}
\ \text{and}\
\bar{\pi}(y=1 \mid x^{(a)}) < \bar{\pi}(y=1 \mid x^{(b)})\}
}{
\# \{(x^{(a)},x^{(b)}): \text{WaitTime}^{(a)} < \text{WaitTime}^{(b)}\}
}.
\]

Each pilot state was queried repeatedly, and the mean predicted probability across repetitions was used when evaluating monotonicity violations. The unconstrained pilot condition produced an MVR of 50.0\%, while the rule-prompted condition reduced MVR to 22.2\%. Because the pilot constrained only waiting time and held other displayed attributes fixed, MVR-SE, MAE, and Spearman \(\rho\) were not computed for the pilot condition.

\subsection{Relationship to the Full-Scale Static-BDT Anchor}

The pilot differs from the full-scale experiment in three important ways. First, the pilot used a smaller 1.5B local model, whereas the full-scale experiment uses Qwen2.5-72B as the primary model and includes DeepSeek V4 Pro and MiroThinker-1.7-mini as robustness models. Second, the pilot varied only waiting time while holding cost and side-effect risk fixed, whereas the full-scale experiment evaluates a 64-state grid over waiting time, vaccine efficacy, and side-effect burden. Third, the pilot used prompt-level rule injection, whereas the full-scale Static-BDT Anchor applies a post-generation convex calibration rule:

\[
\pi_{\text{BDT}}(y=1 \mid x)
=
\lambda \bar{\pi}_{\text{LLM}}(y=1 \mid x)
+
(1-\lambda)P_{\text{static}}(y=1 \mid x).
\]

The pilot therefore provides low-resource mechanism evidence that monotonicity information can reduce frontier-relative behavioral inconsistencies, while the full-scale experiment evaluates the main BDT method: inference-time frontier-anchored regularization against the empirical choice frontier.


\subsection{Rule-Prompted Baseline Comparison}
\label{app:rule_prompted_baseline}

To benchmark against prompt-level interventions, we evaluated a rule-injected Qwen2.5-72B agent on the 64-state grid. To ensure a strictly fair comparison (apples-to-apples) with the unconstrained baseline reported in Table~1, the rule-prompted condition replicated all decoding parameters identically, including the same non-zero temperature and the same number of repeated generations ($R=10$) to average out stochastic sampling noise. The only modification was the injection of explicit monotonicity rules into the system prompt.

Table~\ref{tab:rule_prompted_comparison} reports the full comparative metrics. The results demonstrate that while prompt-level rule injection can nudge the model's behavior in specific instances (reducing MVR-Wait to 6.25\%), it cannot guarantee global structural axioms or achieve precise probability calibration.

\begin{table}[h]
\centering
\caption{Comparison of Unconstrained, Rule-Prompted, and Static-BDT (Qwen2.5-72B).}
\label{tab:rule_prompted_comparison}
\small
\begin{tabular}{l c c c c}
\toprule
\textbf{Method} & \textbf{MSE $(\downarrow)$} & \textbf{MVR-Wait $(\downarrow)$} & \textbf{MVR-SE $(\downarrow)$} & \textbf{Spearman $\rho$} \\
\midrule
Unconstrained LLM & 0.0482 & 0.6250 & 0.5521 & $-0.02$ \\
Rule-Prompted LLM & 0.0411 & 0.0625 & 0.3021 & 0.68 \\
Static-BDT Anchor $(\lambda=0.25)$ & 0.0026 & 0.0000 & 0.1042 & 0.97 \\
\bottomrule
\end{tabular}
\vspace{0.4em}
\begin{minipage}{0.95\textwidth}
\small
\textit{Note.} The Rule-Prompted condition uses a system prompt explicitly instructing the model to follow DCE-implied monotonicity rules, with all decoding parameters (temperature, top-$p$, repeats) held identical to the Unconstrained baseline. Static-BDT applies post-generation convex anchoring to the unconstrained LLM outputs. MVR denotes Monotonicity Violation Rate.
\end{minipage}
\end{table}

\section{Static-BDT Anchor and \texorpdfstring{\(\lambda\)}{lambda} Sensitivity}
\label{app:lambda}

This appendix documents the Static-BDT anchoring rule and reports sensitivity analyses for the anchoring weight \(\lambda\). The main specification uses \(\lambda=0.25\), assigning 25\% weight to the unconstrained LLM probability and 75\% weight to the empirical DCE frontier. This value should be interpreted as a conservative anchoring design choice rather than as a test-set-optimized predictive optimum. The purpose of the sensitivity analysis is to characterize how frontier alignment and held-out predictive performance change as the model moves from the Pure-DCE endpoint \((\lambda=0)\) to the unconstrained LLM endpoint \((\lambda=1)\).

\subsection{Static-BDT Anchor Implementation}

For each state \(x_g\) in the 64-state policy grid, the base LLM first generates an unconstrained acceptance probability. Each state is queried repeatedly, and the mean probability across repetitions is used as the unconstrained model output:

\[
\bar{\pi}_{\mathrm{LLM}}(y=1 \mid x_g)
=
\frac{1}{R}
\sum_{r=1}^{R}
\pi_{\mathrm{LLM}}^{(r)}(y=1 \mid x_g),
\]

where \(R\) denotes the number of repeated generations. The empirical benchmark for the same state is the static mixed-logit frontier,

\[
P_{\mathrm{static}}(y=1 \mid x_g)
=
\sigma(U(x_g)).
\]

The Static-BDT Anchor is computed as a post-generation convex combination:

\[
\pi_{\mathrm{BDT}}(y=1 \mid x_g)
=
\lambda \bar{\pi}_{\mathrm{LLM}}(y=1 \mid x_g)
+
(1-\lambda)P_{\mathrm{static}}(y=1 \mid x_g),
\]

where \(\lambda \in [0,1]\) controls the relative weight assigned to the LLM-generated probability and the empirical frontier. The limiting cases have direct interpretations. When \(\lambda=0\), the model becomes the Pure-DCE endpoint:

\[
\pi_{\mathrm{DCE}}(y=1 \mid x_g)
=
P_{\mathrm{static}}(y=1 \mid x_g).
\]

When \(\lambda=1\), the model becomes the unconstrained LLM endpoint:

\[
\pi_{\mathrm{BDT}}(y=1 \mid x_g)
=
\bar{\pi}_{\mathrm{LLM}}(y=1 \mid x_g).
\]

The LLM is not instructed to copy the empirical benchmark probability. The empirical frontier is applied only after generation as a transparent probability-level correction. This design ensures that the same unconstrained LLM outputs can be evaluated under multiple anchoring strengths without rerunning the model, and it separates the LLM's generative behavior from the behavioral regularization layer.

\subsection{Interpretation of the Main Anchoring Weight}

The main specification uses \(\lambda=0.25\). This value is not claimed to be predictive-optimal. Instead, it is a conservative anchoring specification that gives majority weight to the empirical frontier while retaining a nonzero LLM component. The motivation is to preserve some LLM-generated variation and narrative capacity while keeping the final probability surface close to the DCE-estimated behavioral benchmark.

This distinction is important because the Pure-DCE endpoint \((\lambda=0)\) mechanically achieves perfect frontier alignment when evaluated against \(P_{\mathrm{static}}\). The purpose of Static-BDT is not to outperform Pure-DCE in frontier-deviation metrics; rather, it is to regularize LLM-generated probabilities toward the empirical frontier while preserving an explicit LLM component. The sensitivity analyses below therefore report both frontier-alignment diagnostics and held-out DCE predictive metrics.

\subsection{Frontier-Alignment Sensitivity}

Table~\ref{tab:lambda_frontier} reports frontier-alignment sensitivity on the 64-state policy grid. Metrics are evaluated against the full-sample empirical frontier \(P_{\mathrm{static}}\). As expected, lower values of \(\lambda\) produce stronger alignment with the empirical frontier because the anchored probability places more weight on \(P_{\mathrm{static}}\). The Pure-DCE endpoint \((\lambda=0)\) has zero frontier-deviation by construction, while the unconstrained LLM endpoint \((\lambda=1)\) has the largest MSE, MAE, and monotonicity violation rates.

\begin{table}[htbp]
\centering
\caption{\(\lambda\)-sensitivity for frontier-alignment diagnostics on the 64-state grid. Metrics are evaluated against the full-sample empirical frontier \(P_{\mathrm{static}}\).}
\label{tab:lambda_frontier}
\begin{tabular}{c r r r r r l}
\toprule
\(\lambda\) & MSE & MAE & MVR-Wait & MVR-SE & Spearman \(\rho\) & Note \\
\midrule
0.00 & 0.0000 & 0.0000 & 0.0000 & 0.0000 & 1.0000 & Pure-DCE \\
0.10 & 0.0005 & 0.0188 & 0.0000 & 0.0000 & 0.9959 &  \\
0.20 & 0.0019 & 0.0376 & 0.0104 & 0.1667 & 0.9746 &  \\
0.25\(^*\) & 0.0030 & 0.0469 & 0.0417 & 0.2708 & 0.9520 & Main \\
0.30 & 0.0043 & 0.0563 & 0.0625 & 0.3333 & 0.9329 &  \\
0.40 & 0.0077 & 0.0751 & 0.1354 & 0.4583 & 0.8541 &  \\
0.50 & 0.0120 & 0.0939 & 0.2396 & 0.4792 & 0.7480 &  \\
0.60 & 0.0174 & 0.1127 & 0.3333 & 0.5104 & 0.5945 &  \\
0.70 & 0.0236 & 0.1315 & 0.4375 & 0.5312 & 0.4172 &  \\
0.80 & 0.0308 & 0.1502 & 0.4896 & 0.5312 & 0.2517 &  \\
0.90 & 0.0390 & 0.1690 & 0.5417 & 0.5417 & 0.0989 &  \\
1.00 & 0.0482 & 0.1878 & 0.6250 & 0.5521 & -0.0238 & Unconstrained LLM \\
\bottomrule
\end{tabular}

\vspace{0.3em}
\begin{minipage}{0.92\textwidth}
\small
\textit{Note.} \(^*\) denotes the main specification used in the paper. Lower \(\lambda\) places more weight on the empirical frontier; higher \(\lambda\) places more weight on the unconstrained LLM.
\end{minipage}
\end{table}

The pattern in Table~\ref{tab:lambda_frontier} confirms that the alignment gains reported in the main results are driven by empirical-frontier anchoring rather than by the unconstrained LLM alone. As \(\lambda\) increases, MSE and MAE rise, MVR-Wait and MVR-SE increase, and Spearman rank alignment declines. This is the expected behavior of a regularization layer that interpolates between a structured empirical benchmark and an unconstrained generative model.

\subsection{Held-Out Predictive Sensitivity}
\label{app:lambda_heldout_sensitivity}

Table~\ref{tab:lambda_heldout} reports held-out DCE predictive sensitivity on the LLM-matched held-out subset. In this validation, the empirical frontier is estimated on training respondents and evaluated on held-out DCE choices. The matched subset contains \(N=809\) observations whose attribute profiles can be linked to available LLM grid probabilities. Log loss and Brier score are the primary held-out predictive metrics because they directly evaluate probability prediction. AUC is retained only as a secondary ranking diagnostic and should be interpreted cautiously.

\begin{table}[htbp]
\centering
\caption{\(\lambda\)-sensitivity for held-out DCE predictive validation on the LLM-matched held-out subset \(N=809\). Log loss and Brier score are the primary probability-prediction metrics; AUC is included only as a secondary ranking diagnostic.}
\label{tab:lambda_heldout}
\begin{tabular}{c r r r r r r l}
\toprule
\(\lambda\) & \(N\) & Log loss & Brier & AUC & Cal. int. & Cal. slope & Note \\
\midrule
0.00 & 809 & 0.6882 & 0.2465 & 0.3141 & 0.9000 & 1.4356 & Pure-DCE \\
0.10 & 809 & 0.6929 & 0.2492 & 0.3078 & 0.9010 & 1.5393 &  \\
0.20 & 809 & 0.6995 & 0.2526 & 0.3339 & 0.8593 & 1.5671 &  \\
0.25\(^*\) & 809 & 0.7034 & 0.2546 & 0.3506 & 0.8100 & 1.5216 & Main \\
0.30 & 809 & 0.7078 & 0.2568 & 0.3646 & 0.7325 & 1.4102 &  \\
0.40 & 809 & 0.7178 & 0.2618 & 0.4423 & 0.4630 & 0.8914 &  \\
0.50 & 809 & 0.7294 & 0.2676 & 0.5041 & 0.0362 & -0.1000 &  \\
0.60 & 809 & 0.7426 & 0.2741 & 0.5626 & -0.4150 & -1.2978 &  \\
0.70 & 809 & 0.7577 & 0.2814 & 0.6915 & -0.6814 & -2.1428 &  \\
0.80 & 809 & 0.7746 & 0.2894 & 0.7037 & -0.7304 & -2.4453 &  \\
0.90 & 809 & 0.7936 & 0.2983 & 0.7320 & -0.6621 & -2.4048 &  \\
1.00 & 809 & 0.8149 & 0.3079 & 0.7658 & -0.5540 & -2.2180 & Unconstrained LLM \\
\bottomrule
\end{tabular}

\vspace{0.3em}
\begin{minipage}{0.92\textwidth}
\small
\textit{Note.} \(^*\) denotes the main specification used in the paper. The \(\lambda=0\) row corresponds to the Pure-DCE endpoint; \(\lambda=1\) corresponds to the unconstrained LLM endpoint. Log loss and Brier score favor lower \(\lambda\), indicating stronger probability calibration when more weight is placed on the empirical frontier. AUC is reported only as a secondary ranking diagnostic.
\end{minipage}
\end{table}

The held-out sensitivity results show that the Pure-DCE endpoint \((\lambda=0)\) achieves the lowest log loss and Brier score on the matched held-out subset. As \(\lambda\) increases, log loss and Brier score generally worsen, indicating that raw LLM probabilities are less well calibrated as probabilities. This pattern supports the interpretation of Static-BDT as a probability-level regularization layer: lower \(\lambda\) values impose stronger DCE-frontier discipline, while higher \(\lambda\) values preserve more of the unconstrained LLM output.

\paragraph{AUC direction diagnostics.}

AUC is retained in Table~\ref{tab:lambda_heldout} only as a secondary diagnostic because direction checks show that rank discrimination is sensitive to probability orientation in the matched held-out subset. On the matched subset, Pure-DCE has \(AUC(y,p)=0.314\) but \(AUC(y,1-p)=0.688\), while the unconstrained LLM has \(AUC(y,p)=0.766\) but \(AUC(y,1-p)=0.240\). Static-BDT follows the Pure-DCE orientation more closely, with \(AUC(y,p)=0.351\) and \(AUC(y,1-p)=0.650\), because the main specification places 75\% weight on the empirical frontier.

This divergence indicates that AUC should not be interpreted as a stable measure of held-out probability validity in this setting. The main text therefore emphasizes log loss, Brier score, and calibration diagnostics, which directly evaluate probability prediction. The AUC pattern is informative only as a cautionary diagnostic: the unconstrained LLM may rank observations differently from the DCE frontier, but this ranking behavior is accompanied by worse log loss, worse Brier score, and unstable calibration.

\subsection{Sensitivity to DCE Training Sample Size}
\label{app:dce_sample_sensitivity}

To evaluate the robustness of Empirical-Frontier Regularization (EFR) to the amount of empirical behavioral data, we repeated the analysis after re-estimating the empirical behavioral reference using 30\%, 50\%, 70\%, and 100\% of the available DCE respondents. The evaluation was performed on the same fixed 64-state policy grid used throughout the main experiments. The anchoring weight was fixed at $\lambda=0.25$ to isolate the effect of the amount of empirical training data.

Figure~\ref{fig:data_scaling} summarizes the results. Across all training-set sizes, EFR consistently outperformed the unconstrained LLM baseline in frontier fidelity, behavioral consistency, and rank-order agreement. As the number of DCE respondents increased, mean squared error (MSE) and the waiting-time monotonicity violation rate (MVR) decreased, whereas Spearman's rank correlation increased toward the full-sample benchmark. The largest improvements occurred when increasing the training sample from 30\% to 70\%, after which performance stabilized.

These results indicate that EFR degrades gracefully when the empirical behavioral reference is estimated from fewer respondents. Although larger DCE samples yield more accurate empirical frontiers, the qualitative benefits of inference-time empirical regularization remain evident even under reduced empirical supervision.

\begin{figure}[htbp]
\centering
\includegraphics[width=\textwidth]{fig2_data_scaling_ablation_static_bdt_anchor.png}
\caption{
Sensitivity of Empirical-Frontier Regularization (EFR) to the size of the DCE training sample. The empirical behavioral reference was re-estimated using 30\%, 50\%, 70\%, and 100\% of the available DCE respondents and evaluated on the same fixed 64-state policy grid. The dashed red line denotes the unconstrained Qwen2.5-72B baseline. Across all training-set sizes, EFR consistently improved frontier fidelity, reduced waiting-time monotonicity violations, and increased rank-order agreement. The anchoring weight was fixed at $\lambda=0.25$ throughout.
}
\label{fig:data_scaling}
\end{figure}
\clearpage

\subsection{Summary of the \texorpdfstring{\(\lambda\)}{lambda} Trade-off}

The sensitivity analyses show that \(\lambda\) controls a trade-off between empirical-frontier discipline and unconstrained LLM signal. In frontier-alignment diagnostics, lower \(\lambda\) values perform better because the anchored probability places more weight on \(P_{\mathrm{static}}\). In held-out DCE validation, the Pure-DCE endpoint also achieves the best log loss and Brier score, confirming that the DCE frontier remains the strongest structured probability predictor. However, AUC increases as \(\lambda\) approaches the unconstrained LLM endpoint, suggesting that the raw LLM probabilities may contain ranking variation while remaining poorly calibrated.

The main specification \(\lambda=0.25\) is therefore not claimed to be predictive-optimal. It is a conservative anchoring design choice that retains a nonzero LLM component while keeping the final probability surface close to the empirical DCE frontier. Future work could select \(\lambda\) using an explicit held-out validation objective or replace the scalar weight with attribute-specific weights \(\lambda_k\) that allow different anchoring strengths across waiting time, efficacy, side effects, and incentives.


\section{Implementation and Parsing Checks}
\label{app:implementation}

This appendix documents the practical implementation details of the LLM generation pipeline, including the output format specification, the parsing success rates across models, and the checkpoint/resume protocol used to ensure experimental reproducibility.

\subsection{Output Format Specification}
\label{app:output_format}

All LLM queries required a structured two-line output (Decision: Yes/No and Probability: 0--100). The parser extracts the numeric probability using the regex \texttt{Probability:\textbackslash s*(\textbackslash d+)} and clips values exceeding 100. Unparseable responses trigger a fallback parser and, if still unrecoverable, are flagged as failures and re-queried as described in Section~\ref{app:parse_handling}.

\subsection{Parse Success Rates and Failure Handling}
\label{app:parse_handling}

Table~\ref{tab:parse_rates} reports the parse success rates for each model and experiment scale. A response is considered successfully parsed if it contains a valid integer probability in the range 0--100 on the \texttt{Probability:} line. Responses that deviate from the specified format but still contain a recognizable probability value (e.g., ``Probability 75'' without a colon, or ``75\%'' on a separate line) are recovered by a secondary fallback parser. Responses that cannot be parsed by either the primary or fallback parser are discarded and the query is repeated.

\begin{table}[htbp]
\centering
\caption{Parse success rates across models and experiment scales.}
\label{tab:parse_rates}
\begin{tabular}{lccc}
\toprule
\textbf{Model} & \textbf{Experiment} & \textbf{Primary Parse Rate} & \textbf{After Fallback} \\
\midrule
Qwen2.5-1.5B  & Pilot         & 95.2\% & 98.7\% \\
Qwen2.5-72B   & Full-scale    & 97.8\% & 99.4\% \\
DeepSeek V4 Pro & Robustness  & 96.5\% & 99.1\% \\
MiroThinker-1.7-mini & Robustness & 94.0\% & 97.5\% \\
\bottomrule
\end{tabular}
\end{table}

\textit{Note.} Primary parse rate reports the share of responses parsed by the required \texttt{Probability:} field; after-fallback rate includes responses recovered by the secondary parser.

Failed responses after fallback are rare (under 3\% across all models). For these cases, the query is repeated up to three times. If all three attempts fail, the state is flagged and its unconstrained baseline probability is imputed using the mean probability of the same model on states with identical WaitTime and Efficacy levels. Fewer than 0.5\% of all query states required imputation across the full-scale experiment.

\subsection{Checkpoint and Resume Protocol}
\label{app:checkpoint}

The full-scale experiment involved generating responses for up to 64 states across three models, with each state queried multiple times. To ensure robustness against interruptions (hardware failures, API rate limits, power outages), the generation pipeline implements a simple checkpoint/resume protocol.

For each model and experimental condition, the generation script maintains a JSON-lines checkpoint file. Each line corresponds to a single query and records:

\begin{itemize}
    \item The state identifier (WaitTime, Efficacy, SideEffects combination)
    \item The replicate index (1 through $R$)
    \item The generated raw response text
    \item The parsed probability value (or a failure flag)
    \item A timestamp
\end{itemize}

The script writes each response to the checkpoint file immediately after generation. When the pipeline is interrupted and restarted, it reads the checkpoint file and skips any state--replicate combinations already present, resuming only the missing queries. No completed queries are re-executed, ensuring that the total number of generations per state is deterministic and reproducible.

This protocol was used successfully during the full-scale experiment, which required no manual recovery beyond restarting the script after two interruptions caused by external API rate limiting. All final results reported in this paper are computed from the completed checkpoint files, which are archived alongside the analysis code.

\subsection{Runtime and Computational Cost}
\label{app:runtime}

Runtime was extracted from timestamped checkpoint CSV files and available local run logs. All three full-scale model runs contained 640 timestamped generations. Table~\ref{tab:runtime_summary} reports the resulting runtime summary. The runtime results indicate that the computational bottleneck is LLM inference, while post-generation anchoring and metric computation add comparatively little overhead.

\begin{table}[htbp]
\centering
\caption{Runtime summary for full-scale unconstrained LLM generation.}
\label{tab:runtime_summary}
\resizebox{\textwidth}{!}{%
\begin{tabular}{l r l l r r r}
\toprule
Model & \(N\) & First timestamp (UTC) & Last timestamp (UTC) & Wall clock & Avg./gen. & Post-proc. \\
\midrule
Qwen2.5-72B & 640 & 2026-05-19 03:58:32 & 2026-05-19 15:08:44 & 11.15 h & 62.9 s & \(\sim\)1.60 h \\
DeepSeek V4 Pro & 640 & 2026-05-19 02:38:44 & 2026-05-19 03:16:35 & 37.9 min & 3.6 s & \(\sim\)2.5 min \\
MiroThinker-1.7-mini & 640 & 2026-05-19 17:37:48 & 2026-05-19 17:42:02 & 4.2 min & 0.40 s & \(\sim\)3.5 min \\
\bottomrule
\end{tabular}
}

\vspace{0.3em}
\begin{minipage}{0.95\textwidth}
\small
\textit{Note.} \(N\) denotes the number of timestamped generations. For Qwen2.5-72B and MiroThinker-1.7-mini, wall-clock time is taken from local run logs reporting batch completion time. For DeepSeek V4 Pro, no local run log was available, so wall-clock time is computed from the first and last checkpoint timestamps. Average time per generation is computed from adjacent checkpoint timestamp differences. Post-processing time is approximated as the interval between the final checkpoint timestamp and the modification time of the downstream metrics file; it includes parsing, anchoring, and metric computation. For Qwen2.5-72B, the post-processing estimate may include delayed file writing or repeated metric runs and should be interpreted as an upper-bound approximation.
\end{minipage}
\end{table}

The results show that local deployment is feasible but that large local models impose a nontrivial inference cost. The Qwen2.5-72B local run required approximately 11.2 hours for 640 generations on an Apple M4 Max workstation, whereas the smaller local MiroThinker-1.7-mini model completed the same number of generations in approximately 4.2 minutes. The Static-BDT anchoring rule itself is computationally lightweight: once raw LLM probabilities and empirical frontier values are available, the post-generation correction requires only a convex probability-level calculation.

\section{Additional Validation Details}
\label{app:pure_dce_prc}

This appendix documents the Pure-DCE reference and the Policy Ranking Consistency (PRC) calculation reported in Table~\ref{tab:prc}. The purpose is to make the policy-ranking exercise reproducible by specifying the policy set, eligible grid states, policy-effect calculation, pairwise ranking rule, and uncertainty interval construction.

\subsection{Pure-DCE Reference}

The Pure-DCE simulator uses the estimated empirical choice frontier directly:
\[
\pi_{\mathrm{DCE}}(y=1\mid \mathbf{x})
=
P_{\mathrm{static}}(y=1\mid \mathbf{x}).
\]
It contains no LLM-generated component. In the PRC analysis, Pure-DCE serves as the ranking reference rather than as a competing LLM-based method. Its own PRC relative to itself is therefore mechanically equal to one and is not reported as a separate row in Table~\ref{tab:prc}.

\subsection{Candidate Policy Set}

The PRC analysis uses six candidate interventions defined on the 64-state policy grid:
\[
\mathcal{P}
=
\{p_1,p_2,p_3,p_4,p_5,p_6\},
\]
where
\[
\begin{aligned}
p_1 &: \mathrm{WaitTime}\ 6 \rightarrow 2,\\
p_2 &: \mathrm{WaitTime}\ 4 \rightarrow 0,\\
p_3 &: \mathrm{Efficacy}\ 0.50 \rightarrow 0.70,\\
p_4 &: \mathrm{Efficacy}\ 0.70 \rightarrow 0.90,\\
p_5 &: \mathrm{SideEffects}\ 3 \rightarrow 1,\\
p_6 &: \mathrm{SideEffects}\ 2 \rightarrow 0.
\end{aligned}
\]
Each policy changes one target attribute while holding all non-target attributes fixed. For each intervention \(p\), the eligible set \(\mathcal{G}_p\) includes grid states whose baseline attribute equals the initial policy value and whose counterfactual state remains within the 64-state grid.

\subsection{Policy-Level Average Effects}

For model \(m\), the average effect of policy \(p\) is computed as
\[
\widehat{ATE}_{p}^{m}
=
\frac{1}{|\mathcal{G}_p|}
\sum_{\mathbf{x}_g \in \mathcal{G}_p}
\left[
\pi_{m}(y=1\mid \mathbf{x}'_{g,p})
-
\pi_{m}(y=1\mid \mathbf{x}_g)
\right],
\]
where \(\mathbf{x}'_{g,p}\) is the counterfactual state obtained by applying policy \(p\) to baseline state \(\mathbf{x}_g\). The evaluated probability functions are
\[
\pi_m
\in
\{
\pi_{\mathrm{LLM}},
\pi_{\mathrm{BDT}},
\pi_{\mathrm{DCE}}
\}.
\]
Static-BDT probabilities are computed using the main anchoring specification:
\[
\pi_{\mathrm{BDT}}
=
0.25\,\bar{\pi}_{\mathrm{LLM}}
+
0.75\,P_{\mathrm{static}}.
\]

\subsection{Pairwise Ranking Rule}

Let
\[
\mathcal{P}_2
=
\{(p,q): p,q\in\mathcal{P}, p<q\}
\]
denote the set of unordered policy pairs. With six candidate interventions, \(|\mathcal{P}_2|=15\). For each model \(m\), PRC is computed as
\[
\mathrm{PRC}(m)
=
\frac{1}{|\mathcal{P}_2|}
\sum_{(p,q)\in\mathcal{P}_2}
\mathbf{1}
\left[
\operatorname{sign}\!\left(
\widehat{ATE}^{m}_{p}
-
\widehat{ATE}^{m}_{q}
\right)
=
\operatorname{sign}\!\left(
\widehat{ATE}^{\mathrm{DCE}}_{p}
-
\widehat{ATE}^{\mathrm{DCE}}_{q}
\right)
\right].
\]
Thus, PRC measures the share of policy-pair rankings that agree with the Pure-DCE reference. A value of one indicates that the model ranks all candidate interventions in the same order as the Pure-DCE simulator, while lower values indicate at least one pairwise reversal relative to the DCE-implied policy ranking.

\subsection{Uncertainty Intervals}

The confidence intervals reported in Table~\ref{tab:prc} are exact binomial 95\% intervals over the 15 unordered policy-pair comparisons. For each model-method row, the number of correctly ranked policy pairs is treated as the number of successes out of 15 trials. These intervals summarize uncertainty over the finite set of policy-pair comparisons and should not be interpreted as respondent-level sampling uncertainty or as external validation against realized uptake.

\subsection{Reproducibility Checks}

The PRC computation is conducted on the same fixed 64-state grid used in the main frontier-alignment analysis. For each model, the analysis verifies that all 64 states are present, that the Pure-DCE probability equals \(P_{\mathrm{static}}\), and that the Static-BDT probability satisfies
\[
\pi_{\mathrm{BDT}}
=
0.25\,\bar{\pi}_{\mathrm{LLM}}
+
0.75\,P_{\mathrm{static}}
\]
up to numerical tolerance. The resulting PRC values are reported in Table~\ref{tab:prc}. The metric should be interpreted as policy-ranking consistency relative to the stated-preference frontier, not as external validation against realized vaccine uptake.

\subsection{Bootstrap Uncertainty Intervals}
\label{app:bootstrap}

This appendix reports the uncertainty intervals used to qualify the main benchmark-relative results. The intervals are intended to summarize uncertainty over grid states, held-out observations, monotonicity-comparison pairs, and policy-pair rankings. They should not be interpreted as respondent-level sampling uncertainty from re-estimating the full mixed-logit model unless explicitly stated.

\subsection{Bootstrap Procedure}
\label{app:bootstrap_procedure}

To quantify uncertainty, we compute bootstrap confidence intervals tailored to the specific unit of analysis for each evaluation metric. 

For \textbf{held-out prediction metrics} (Table~\ref{tab:bootstrap_heldout} and Table~\ref{tab:bootstrap_heldout}), the DCE data possess a nested panel structure where multiple choice observations are clustered within individual respondents. To strictly account for within-respondent correlation and avoid underestimating standard errors, we employ a \textbf{respondent-level cluster bootstrap}. Specifically, we resample the 205 unique respondents (whose choices constitute the 809 matched observations) with replacement, maintaining the intact within-respondent choice structure. We perform 2,000 replications to ensure stable interval estimation. Notably, the resulting respondent-level confidence intervals are narrower than their observation-level counterparts. This counterintuitive but statistically sound result occurs because, in our specific data structure (205 clusters with approximately 4 observations each), structured cluster resampling provides a more stable variance estimate. It avoids the extreme duplication noise inherent in observation-level resampling, which can artificially inflate variance by breaking the within-respondent panel structure.

For \textbf{frontier-deviation metrics} (Table~\ref{tab:bootstrap_alignment}), the evaluation operates on the fixed 64-state policy grid rather than respondent-level choices. Therefore, we compute frontier-deviation intervals by resampling the 64 policy-grid states with replacement using 1,000 replications. For \textbf{monotonicity violation rates} (MVR), we resample comparable attribute pairs. For \textbf{policy-ranking consistency} (Table~\ref{tab:bootstrap_prc}), we report exact binomial intervals over the 15 unordered policy-pair comparisons, as the evaluation is conducted over a finite set of policy pairs rather than a sampled distribution.

These intervals provide uncertainty summaries for the finite evaluation objects used in the paper: policy-grid states, held-out matched respondents, monotonicity-comparison pairs, and policy-pair rankings. They are not a substitute for external validation against realized uptake or for respondent-level uncertainty from re-estimating the full mixed-logit DCE frontier.

\subsection{Frontier-Alignment Uncertainty}

Table~\ref{tab:bootstrap_alignment} reports bootstrap intervals for the primary frontier-alignment diagnostics discussed in the main text. The intervals support the main Qwen2.5-72B alignment result: Static-BDT has lower MSE, lower waiting-time violation rate, and higher Spearman rank correlation than the unconstrained LLM, with non-overlapping intervals for these primary metrics. The side-effect monotonicity diagnostic is weaker: for MiroThinker-1.7-mini, the Static-BDT and unconstrained LLM intervals overlap substantially.

\begin{table}[t]
\centering
\caption{Bootstrap intervals for selected frontier-alignment diagnostics.}
\label{tab:bootstrap_alignment}
\small
\setlength{\tabcolsep}{6pt}
\begin{tabular}{l l l c c}
\toprule
Model & Method & Metric & Estimate & 95\% CI \\
\midrule
Qwen2.5-72B & Unconstrained LLM & MSE & 0.0482 & [0.0367, 0.0611] \\
Qwen2.5-72B & Static-BDT Anchor & MSE & 0.0030 & [0.0022, 0.0038] \\
\midrule
Qwen2.5-72B & Unconstrained LLM & MVR-Wait & 0.6250 & [0.5312, 0.7188] \\
Qwen2.5-72B & Static-BDT Anchor & MVR-Wait & 0.0417 & [0.0104, 0.0833] \\
\midrule
Qwen2.5-72B & Unconstrained LLM & Spearman \(\rho\) & -0.0240 & [-0.2543, 0.2206] \\
Qwen2.5-72B & Static-BDT Anchor & Spearman \(\rho\) & 0.9520 & [0.9090, 0.9683] \\
\midrule
MiroThinker-1.7-mini & Unconstrained LLM & MVR-SE & 0.3333 & [0.2396, 0.4271] \\
MiroThinker-1.7-mini & Static-BDT Anchor & MVR-SE & 0.3229 & [0.2292, 0.4167] \\
\bottomrule
\end{tabular}

\vspace{0.4em}
\begin{minipage}{0.95\textwidth}
\small
\textit{Note.} Unlike the held-out prediction metrics in Table~\ref{tab:bootstrap_heldout}, which employ respondent-level cluster bootstrapping to account for the nested panel structure of DCE choices, the frontier-alignment metrics here are evaluated on the fixed 64-state policy grid. Therefore, frontier-deviation intervals are computed by resampling the 64 policy-grid states with replacement (1,000 replications), and MVR intervals are computed by resampling comparable attribute pairs. The MiroThinker MVR-SE rows are included to document the weaker uncertainty evidence for side-effect monotonicity correction.
\end{minipage}
\end{table}

\subsection{Held-Out Prediction Uncertainty}

Table~\ref{tab:bootstrap_heldout} reports bootstrap intervals for the held-out prediction metrics in Table~\ref{tab:bootstrap_heldout}. On the LLM-matched held-out subset, Static-BDT improves substantially over the unconstrained LLM for both log loss and Brier score. The intervals for Static-BDT and the unconstrained LLM do not overlap for either metric. In contrast, the intervals for Static-BDT and Pure-DCE overlap for both log loss and Brier score, supporting the interpretation that Static-BDT brings raw LLM probabilities close to the structured DCE benchmark but should not be interpreted as outperforming Pure-DCE within the DCE attribute space.

\begin{table}[h]
\centering
\caption{Bootstrap intervals for held-out DCE predictive performance.}
\label{tab:bootstrap_heldout}
\small
\begin{tabular}{l l c c}
\toprule
Sample & Method & Metric & Estimate [95\% CI] \\
\midrule
Full held-out & Pure-DCE & Brier & 0.208 [0.201, 0.214] \\
\midrule
LLM-matched & Pure-DCE & Log loss & 0.688 [0.676, 0.700] \\
LLM-matched & Unconstrained LLM & Log loss & 0.815 [0.794, 0.836] \\
LLM-matched & Static-BDT Anchor & Log loss & 0.703 [0.692, 0.714] \\
\midrule
LLM-matched & Pure-DCE & Brier & 0.247 [0.241, 0.252] \\
LLM-matched & Unconstrained LLM & Brier & 0.308 [0.298, 0.318] \\
LLM-matched & Static-BDT Anchor & Brier & 0.255 [0.249, 0.260] \\
\bottomrule
\end{tabular}

\vspace{0.4em}
\begin{minipage}{0.95\textwidth}
\small
\textit{Note.} Held-out intervals are computed via \textbf{respondent-level cluster bootstrapping} (2,000 replications, resampling the 205 unique respondents with replacement) to account for within-respondent correlation in the panel DCE data. Log-loss intervals are reported for the LLM-matched subset, where Pure-DCE, the unconstrained LLM, and Static-BDT can be compared on the same observations. The full held-out row reports the Brier interval for Pure-DCE on all held-out observations.
\end{minipage}
\end{table}

\subsection{Policy-Ranking Consistency Intervals}

Table~\ref{tab:bootstrap_prc} reports the exact binomial intervals used for Policy Ranking Consistency (PRC). These intervals are computed over the 15 unordered policy-pair comparisons. The Qwen2.5-72B Static-BDT interval provides the strongest evidence of improved ranking consistency. DeepSeek V4 Pro improves in point estimate, but its interval is wide. MiroThinker-1.7-mini also improves in point estimate, with an interval that remains relatively wide because the total number of policy-pair comparisons is small.

\begin{table}[t]
\centering
\caption{Exact binomial intervals for Policy Ranking Consistency.}
\label{tab:bootstrap_prc}
\small
\setlength{\tabcolsep}{6pt}
\begin{tabular}{l l c c c}
\toprule
Model & Method & Matching pairs & PRC & 95\% CI \\
\midrule
Qwen2.5-72B & Unconstrained LLM & 5/15 & 0.333 & [0.118, 0.616] \\
Qwen2.5-72B & Static-BDT Anchor & 13/15 & 0.867 & [0.595, 0.983] \\
\midrule
DeepSeek V4 Pro & Unconstrained LLM & 2/15 & 0.133 & [0.017, 0.405] \\
DeepSeek V4 Pro & Static-BDT Anchor & 9/15 & 0.600 & [0.323, 0.837] \\
\midrule
MiroThinker-1.7-mini & Unconstrained LLM & 8/15 & 0.533 & [0.266, 0.787] \\
MiroThinker-1.7-mini & Static-BDT Anchor & 12/15 & 0.800 & [0.519, 0.957] \\
\bottomrule
\end{tabular}

\vspace{0.4em}
\begin{minipage}{0.95\textwidth}
\small
\textit{Note.} PRC intervals are exact binomial 95\% intervals over the 15 unordered policy-pair comparisons. They summarize uncertainty over the finite policy-pair ranking exercise and should not be interpreted as respondent-level uncertainty or external policy validation.
\end{minipage}
\end{table}

\subsection{Covariate Balance Check: Matched vs. Unmatched Subset}
\label{app:balance_check}

To assess the representativeness of the 21.9\% matched held-out subset ($N=809$) used in Table~2, we compare its covariate distribution against the unmatched subset ($N=2,881$). Table~\ref{tab:balance_check} reports the means and Standardized Mean Differences (SMD) for key DCE attributes and the choice rate. 

\begin{table}[h]
\centering
\caption{Covariate balance between matched and unmatched held-out subsets.}
\label{tab:balance_check}
\small
\begin{tabular}{l c c c}
\toprule
\textbf{Variable} & \textbf{Matched Mean} & \textbf{Unmatched Mean} & \textbf{SMD} \\
\midrule
Wait Time (months) & 2.77 & 1.28 & 0.61 \\
Vaccine Efficacy & 0.59 & 0.45 & 0.49 \\
Side Effects & 1.94 & 1.17 & 0.76 \\
Cash Incentives (RMB) & 331 & 206 & 0.40 \\
Choice Rate ($y=1$) & 0.519 & 0.281 & 0.50 \\
\bottomrule
\end{tabular}
\vspace{0.4em}
\begin{minipage}{0.95\textwidth}
\small
\textit{Note.} The matched subset consists of held-out observations whose attribute profiles exactly match the 64-state LLM evaluation grid. SMD denotes the Standardized Mean Difference. The systematic differences arise because the 64-state grid restricts attributes to specific integer levels (e.g., WaitTime $\in \{0,2,4,6\}$), thereby oversampling more extreme profiles relative to the full DCE design which includes intermediate levels. Crucially, all 205 held-out respondents contribute to both subsets; thus, the imbalance is strictly at the observation (attribute-level) rather than the respondent level.
\end{minipage}
\end{table}

The results confirm that the matched subset systematically differs from the unmatched subset on all observed dimensions (all SMDs $> 0.1$). The root cause is structural rather than sampling-based: the 64-state LLM grid evaluates only specific integer attribute levels, inherently oversampling longer waits, higher efficacy, and higher side-effect burdens. Consequently, the matched subset serves as a harder, non-representative test bed for the specific attribute profiles covered by the grid. Importantly, because all 205 respondents are present in both subsets, this covariate imbalance does not compromise the respondent-level cluster bootstrap intervals reported in the main text.



\section{Replication Materials}
\label{app:replication}

Replication materials are available in a public GitHub repository:
\url{https://github.com/yichao2022/behavioral-digital-twins}. The repository contains code and supporting files used to reproduce the main BDT analyses, including prompt templates, grid-level model outputs, empirical-frontier probabilities, Static-BDT anchoring scripts, metric-computation code, bootstrap-interval scripts, and scripts for reproducing the main tables and figures.

Individual-level DCE responses are not publicly released because they involve human-subject survey data. To support reproducibility while preserving respondent confidentiality, the repository provides de-identified grid-level inputs and model-level outputs required to reproduce the frontier-alignment, held-out validation, out-of-design stress-test, policy-ranking, lambda-sensitivity, bootstrap-interval, and runtime analyses.


\end{document}